\documentclass[11pt,a4paper,openany]{report}
\usepackage[utf8]{inputenc}
\usepackage[T1]{fontenc}
\usepackage{lmodern}
\usepackage[margin=2.4cm]{geometry}
\usepackage{microtype}
\usepackage{eurosym}
\usepackage{amsmath,amssymb,amsthm,mathtools}
\usepackage{booktabs,array,multirow}
\usepackage{enumitem}
\usepackage{xcolor}
\usepackage[most]{tcolorbox}
\usepackage{tikz}
\usetikzlibrary{arrows.meta,positioning,calc,patterns}
\usepackage{pgfplots}
\pgfplotsset{compat=1.17}
\usepackage[hidelinks]{hyperref}
\usepackage{fancyhdr}
\usepackage{titlesec}

% --- lingua (senza babel-italian) ---
\renewcommand{\contentsname}{Indice}
\renewcommand{\chaptername}{Capitolo}
\renewcommand{\figurename}{Figura}
\renewcommand{\tablename}{Tabella}
\setlength{\parindent}{0pt}
\setlength{\parskip}{0.45em}
\sloppy

\titleformat{\chapter}[display]{\normalfont\huge\bfseries\color{blue!50!black}}{\large\chaptertitlename\ \thechapter}{8pt}{\Huge}
\titlespacing*{\chapter}{0pt}{-20pt}{20pt}

\pagestyle{fancy}
\fancyhf{}
\fancyhead[L]{\small\textsc{Fondamenti di Ricerca Operativa}}
\fancyhead[R]{\small\nouppercase{\leftmark}}
\fancyfoot[C]{\thepage}

% --- box ---
\newtcolorbox[auto counter,number within=chapter]{defbox}[1][]{enhanced,breakable,colback=blue!4,colframe=blue!50!black,fonttitle=\bfseries,title={Definizione~\thetcbcounter\ifx&#1&\else\ (#1)\fi}}
\newtcolorbox{esbox}[1][]{enhanced,breakable,colback=green!3,colframe=green!40!black,fonttitle=\bfseries,title={Esempio\ifx&#1&\else: #1\fi}}
\newtcolorbox{oss}[1][]{enhanced,breakable,colback=gray!6,colframe=gray!60,fonttitle=\bfseries,coltitle=black,title={Osservazione\ifx&#1&\else: #1\fi}}
\newtcolorbox{esame}{enhanced,breakable,colback=orange!6,colframe=orange!70!black,fonttitle=\bfseries,title={Per l'esame}}
\newtcolorbox{extra}[1][]{enhanced,breakable,colback=violet!4,colframe=violet!60!black,fonttitle=\bfseries,title={Approfondimento\ifx&#1&\else: #1\fi}}
\newtcolorbox{modello}[1][]{enhanced,breakable,colback=white,colframe=black!70,boxrule=0.6pt,fonttitle=\bfseries,coltitle=black,colbacktitle=black!8,title={#1}}

\newtheorem{prop}{Proposizione}[chapter]
\theoremstyle{remark}
\newtheorem*{dimo}{Dimostrazione}

% --- macro ---
\newcommand{\R}{\mathbb{R}}
\newcommand{\Z}{\mathbb{Z}}
\newcommand{\N}{\mathbb{N}}
\newcommand{\st}{\text{s.t.}}
\newcommand{\bx}{\mathbf{x}}
\newcommand{\NP}{\mathcal{NP}}
\newcommand{\Pc}{\mathcal{P}}
\newcommand{\red}{\propto}

\begin{document}

\begin{titlepage}
\centering
\vspace*{3cm}
{\Huge\bfseries Fondamenti di Ricerca Operativa\par}
\vspace{0.6cm}
{\Large Appunti del corso -- A.A.\ 2026/27\par}
\vspace{0.4cm}
{\large Prof.\ Marta Pascoal (Politecnico di Milano, DEIB)\par}
\vspace{1.5cm}
{\large Lezioni 1--5 (15--25 settembre 2026) + PL dalle slide 04\par}
\vspace{0.3cm}
\begin{minipage}{0.75\textwidth}\centering
Introduzione alla RO $\cdot$ Modelli di Programmazione Lineare (intera e mista) $\cdot$ Variabili libere e discrete $\cdot$ Cenni di complessità computazionale $\cdot$ PL: forme, geometria e teorema fondamentale $\cdot$ Quesiti V/F di allenamento $\cdot$ Esercitazione 1 (modelli) con soluzioni
\end{minipage}
\vfill
{\small Redatti a partire dalle slide del corso (adattate da G.\ Carello), dagli appunti di P.\ Belotti (versione 2025) e dall'Esercitazione n.\,1. I box viola \emph{Approfondimento} contengono materiale non presente nelle slide di quest'anno ma utile per capire (o che servirà più avanti). I risultati numerici degli esempi sono stati verificati con un solutore.\par}
\end{titlepage}

\tableofcontents

\chapter*{Informazioni sul corso e sull'esame}
\addcontentsline{toc}{chapter}{Informazioni sul corso e sull'esame}
\markboth{Informazioni}{Informazioni}

\textbf{Docente:} Marta Pascoal (DEIB, edificio 20, II piano), \texttt{marta.brazpascoal@polimi.it}, ricevimento su appuntamento.

\textbf{Orari:} martedì 14:15--16:15 (aula 3.0.2, lezione); venerdì 10:15--12:15 (aula 25.S.2, lezione/esercitazione); venerdì 13:15--15:15 (aula 3.0.2, laboratorio/recuperi). Tutto registrato e in streaming.

\textbf{Esame.}
\begin{itemize}[nosep]
  \item Scritto (in presenza), max 30 punti: esercizi, \emph{domande teoriche} e conoscenza del \emph{linguaggio di modellazione} visto in laboratorio.
  \item Niente libri, appunti, calcolatrici o dispositivi (possederli comporta penalità o annullamento).
  \item Punteggio $\le 10$ $\Rightarrow$ respinto. Nessuna prova in itinere. Eventuale orale a discrezione del docente, \emph{non integrativo} (può anche abbassare il voto).
  \item \textbf{Quesiti brevi}: 4 quiz da 15 minuti, all'inizio della lezione del martedì, solo in presenza; valgono in tutto 3 punti (si scarta il peggiore); validi solo per l'A.A.\ in corso.
\end{itemize}

\textbf{Quesiti brevi: come funzionano} (regole pubblicate a ottobre).
\begin{itemize}[nosep]
  \item Si svolgono su Moodle \textbf{dallo smartphone} (link su WeBeep all'inizio della prova): aggiorna prima il sistema operativo. Durata 15 minuti; niente materiale, niente calcolatrice, niente domande al docente.
  \item Formato \textbf{Vero/Falso}: ogni domanda ha più affermazioni (es.\ Q1.1--Q1.4). Le stesse risposte (``V''/``F'') vanno ricopiate sul foglio cartaceo con nome, cognome e codice persona.
  \item Ogni quesito vale tra 0 e 1 punto, frazionario in base alle risposte corrette, senza arrotondamento; si scarta il peggiore dei 4. I punti si sommano al punteggio dello scritto \emph{solo se lo scritto è sufficiente}; si arrotonda solo il totale.
  \item Per allenarti: Cap.~\ref{cap:quiz} (affermazioni V/F con soluzioni).
\end{itemize}

\textbf{Date:} quesiti brevi il 6/10, 3/11, 24/11, 15/12 (ore 14:15). Laboratori il 23/10, 6/11, 27/11, 11/12 (ven.\ 13:15). Lezioni sospese il 29/9 e il 2/10. \textbf{Primo quesito breve: martedì 6/10} (presumibilmente sugli argomenti dei Capitoli 1--3).

\textbf{Programma del corso}
\begin{enumerate}[nosep]
  \item Modellazione in Programmazione Lineare (PL): tipi di variabili, vincoli, funzioni obiettivo. \hfill \emph{(questi appunti)}
  \item Cenni di complessità computazionale. \hfill \emph{(questi appunti)}
  \item Proprietà della PL e metodo del simplesso.
  \item Metodi per la Programmazione Lineare Intera (PLI).
  \item Problemi di ottimizzazione su grafo.
  \item Laboratorio: linguaggio di modellazione e solutori.
\end{enumerate}

\textbf{Mappa delle lezioni coperte}

\begin{tabular}{@{}llp{0.78\textwidth}@{}}
\toprule
Lezione & Data & Contenuto \\ \midrule
1 & 15/9 & Introduzione, esempio delle sale operatorie (\S\ref{sec:sale}) \\
2 & 18/9 & Soluzioni ammissibili/ottime, forma parametrica e compatta; coltivatore, dieta (\S\ref{sec:def}, \S\ref{sec:colt}, \S\ref{sec:dieta}) \\
3 & 22/9 & Investimento, zaino, big-$M$, vincoli logici, blending, localizzazione bottleneck (\S\ref{sec:zaino}--\S\ref{sec:bottleneck}) \\
4 & 25/9 & Variabili libere e discrete (\S\ref{sec:libere}, \S\ref{sec:discrete}); complessità (Cap.~\ref{cap:compl}) \\
5 & 25/9 & Esercitazione 1: modelli, in aula esercizi 1, 6, 8, 9 (Cap.~\ref{cap:es}) \\
-- & dal 6/10 & Programmazione Lineare: forme, geometria, teorema fondamentale (Cap.~\ref{cap:pl}, scritto dalle slide 04 prima della lezione) \\
\bottomrule
\end{tabular}

\chapter{Introduzione alla Ricerca Operativa}

\section{Che cos'è la Ricerca Operativa}

La \textbf{Ricerca Operativa} (RO, in inglese \emph{Operations Research} o \emph{Management Science}) è la disciplina che fornisce metodi quantitativi e tecniche matematiche (ottimizzazione, teoria dei giochi, simulazione) per risolvere \emph{problemi decisionali complessi} in cui si devono:
\begin{itemize}[nosep]
  \item gestire \textbf{risorse limitate} (non tutte le scelte sono accettabili);
  \item tenere conto di un \textbf{obiettivo} per valutare la bontà delle diverse soluzioni;
  \item \textbf{selezionare} una tra le molte alternative così da ottimizzare l'obiettivo.
\end{itemize}
È una ``scienza di supporto alle decisioni'', all'incrocio tra matematica applicata, informatica, economia e ingegneria gestionale. Nasce durante la II Guerra Mondiale (``research on operations'': posizionamento dei radar, logistica militare) e poi si diffonde in ambito industriale e sociale.

Alcune date da conoscere (non serve impararle a memoria, ma danno il quadro): metodo del simplesso (G.\ Dantzig, 1947); algoritmo di Bellman--Ford per cammini minimi (1956--58); algoritmo di Ford--Fulkerson per il flusso massimo (1956); piani di taglio di Gomory (metà anni '50); algoritmo di Duan et al.\ (2025) per il cammino minimo, che batte Dijkstra in certi regimi.

Esempi tipici: turnazione del personale ospedaliero, percorso più breve su una rete stradale, gestione di un portafoglio di investimenti, pianificazione delle sale operatorie.

\subsection{Il processo di ottimizzazione}

\begin{center}
\begin{tikzpicture}[node distance=0.9cm and 0.7cm, box/.style={draw,rounded corners,fill=blue!6,minimum width=2.0cm,minimum height=0.8cm,font=\small}, >=Stealth]
\node[box] (p) {Problema};
\node[box,right=of p] (m) {Modello};
\node[box,right=of m] (a) {Metodo di soluzione};
\node[box,right=of a] (s) {Soluzione};
\node[box,below=of s] (i) {Implementazione};
\node[box,below=of m] (f) {Feedback esperti};
\draw[->] (p)--(m); \draw[->] (m)--(a); \draw[->] (a)--(s); \draw[->] (s)--(i);
\draw[->] (i)--(f); \draw[->] (f)--(p);
\end{tikzpicture}
\end{center}
\begin{enumerate}[nosep]
  \item si definisce il problema da modellare;
  \item lo si rappresenta con un \textbf{modello matematico} (formulazione in Programmazione Matematica);
  \item lo si risolve con tecniche di ottimizzazione;
  \item si analizza la soluzione con gli esperti del dominio: risponde alle loro esigenze? Se no, si torna al modello (vincoli dimenticati, obiettivo sbagliato, dati imprecisi\ldots).
\end{enumerate}

\section{Un esempio applicativo: le sale operatorie}\label{sec:sale}

\textbf{Problema.} Ci sono 10 pazienti in lista d'attesa, 2 giorni, 2 sale operatorie aperte dalle 7 alle 15 (8 ore). Si hanno quindi $2\times 2 = 4$ \emph{sedute operatorie} (blocchi giorno--sala) da 8 ore. Le durate degli interventi sono:

\begin{center}
\begin{tabular}{@{}l*{10}{c}@{}}
\toprule
Paziente & 1 & 2 & 3 & 4 & 5 & 6 & 7 & 8 & 9 & 10\\
Durata (ore) & 3 & 3 & 3.5 & 4 & 4.5 & 3 & 5 & 6 & 5.5 & 6\\
\bottomrule
\end{tabular}
\end{center}
\emph{Quanti pazienti riusciamo a operare al massimo? E come dimostriamo che è il miglior risultato possibile?}

Provando ``a mano'' si trova facilmente una soluzione con 7 pazienti, ma per \emph{dimostrare} che non si può fare di meglio serve un ragionamento (o un modello e un solutore). Il modello risponde a tre domande:
\begin{center}
\begin{tabular}{@{}lcl@{}}
\toprule
\textbf{Problema} & & \textbf{Modello}\\ \midrule
Descrizione dei dati & $\to$ & insiemi e parametri\\
Cosa dobbiamo decidere? & $\to$ & variabili decisionali\\
Che requisiti rispettare? & $\to$ & vincoli\\
Come valutiamo le scelte? & $\to$ & funzione obiettivo\\
\bottomrule
\end{tabular}
\end{center}

\paragraph{Forma ``esplicita''.} Decidiamo a quale seduta assegnare ogni paziente: $x_{\text{paz}}^{\text{giorno,sala}}\in\{0,1\}$ vale 1 se il paziente è operato in quel giorno e in quella sala. Vincoli: ogni paziente al più una volta,
\[ x_{1}^{g1,s1}+x_{1}^{g1,s2}+x_{1}^{g2,s1}+x_{1}^{g2,s2}\le 1 \quad(\text{e così per ogni paziente}),\]
e non superare le 8 ore di ogni seduta,
\[ 3x_1^{g1,s1}+3x_2^{g1,s1}+3.5x_3^{g1,s1}+\dots+6x_{10}^{g1,s1}\le 8 \quad(\text{e così per ogni seduta}).\]
Obiettivo: massimizzare il numero di pazienti operati, cioè la somma di tutte le $x$.

\paragraph{Forma parametrica.} Scrivere i vincoli uno a uno è scomodo e lega il modello a \emph{quella} istanza. Generalizziamo:
\begin{itemize}[nosep]
  \item $I$: insieme dei pazienti; $J$: insieme delle sedute operatorie;
  \item $t_i$: durata prevista dell'intervento del paziente $i\in I$; $\gamma_j$: durata della seduta $j\in J$;
  \item $x_{ij}\in\{0,1\}$: vale 1 se il paziente $i$ è assegnato alla seduta $j$.
\end{itemize}
\begin{modello}{Modello delle sale operatorie (forma parametrica)}
\begin{align*}
\max\ & \sum_{i\in I}\sum_{j\in J} x_{ij} && \text{(n.\ pazienti operati)}\\
\st\ & \sum_{j\in J} x_{ij}\le 1 && \forall i\in I \quad\text{(ogni paziente al più una volta)}\\
& \sum_{i\in I} t_i x_{ij}\le \gamma_j && \forall j\in J \quad\text{(capacità di ogni seduta)}\\
& x_{ij}\in\{0,1\} && \forall i\in I,\ j\in J
\end{align*}
\end{modello}
Nota: $\sum_{j} x_{ij}$ vale 1 se e solo se $i$ è operato, quindi la doppia sommatoria conta i pazienti operati. Il modello è di \textbf{Programmazione Lineare Intera} (binaria).

\begin{esbox}[perché la risposta è 7]
Una soluzione con 7 pazienti: seduta A $=\{5\}$ (4.5\,h), B $=\{1,7\}$ (8\,h), C $=\{2,4\}$ (7\,h), D $=\{3,6\}$ (6.5\,h). \\
\textbf{Non si può fare 8.} I pazienti con durata $\ge 5.5$ (cioè 8, 9, 10) non possono condividere la seduta con nessun altro, perché resterebbero al più $2.5$\,h e l'intervento più breve dura 3\,h. I pazienti ``piccoli'' (durata $\le 5$) sono solo 7: $\{1,2,3,4,5,6,7\}$, con durata totale $3+3+3.5+4+4.5+3+5=26$.
\begin{itemize}[nosep]
\item Se tra gli 8 operati c'è esattamente un paziente ``grande'', occupa una seduta da solo; gli altri 7 sono tutti i piccoli (26 ore) e dovrebbero stare in 3 sedute ($24$ ore): impossibile.
\item Con 2 grandi: 6 piccoli in 2 sedute (16 ore), ma i 6 piccoli più brevi durano già $21$ ore: impossibile. Con 3 grandi: 5 piccoli in una seduta: impossibile.
\end{itemize}
Quindi l'ottimo vale 7. Questo tipo di argomento (trovare una soluzione \emph{e} un limite superiore che coincidono) è esattamente ciò che fanno, in modo sistematico, gli algoritmi che vedremo.
\end{esbox}

\section{Definizioni fondamentali}\label{sec:def}

Un \textbf{problema di ottimizzazione} è
\[ \min\{f(x): x\in X\}\quad\text{oppure}\quad \max\{f(x):x\in X\},\]
dove $X$ è l'insieme delle soluzioni ammissibili e $f$ la funzione obiettivo. ``s.t.'' (\emph{subject to}) separa l'obiettivo dai vincoli. Nel corso consideriamo spesso la forma
\[ (P)\qquad \max\ c^T x\quad \st\ Ax\le b,\ x\ge 0. \]

\begin{defbox}[soluzione e regione ammissibile]
Assegnando un valore alle variabili si ottiene una \emph{soluzione}. $\bar x$ è una \textbf{soluzione ammissibile} di $(P)$ se soddisfa tutti i vincoli: $A\bar x\le b$ e $\bar x\ge 0$. L'insieme delle soluzioni ammissibili $X=\{x: Ax\le b,\ x\ge 0\}$ è la \textbf{regione ammissibile}.
\end{defbox}

\begin{defbox}[soluzione ottima]
$x^*\in X$ è una \textbf{soluzione (globalmente) ottima} per un problema di massimo se $c^Tx^*\ge c^Tx$ per ogni $x\in X$; per un problema di minimo se $c^Tx^*\le c^Tx$ per ogni $x\in X$. Il valore $c^Tx^*$ è il \textbf{valore ottimo}.
\end{defbox}

Nell'esempio delle sale operatorie
$X=\big\{x: \sum_{j}x_{ij}\le 1\ \forall i;\ \sum_i t_ix_{ij}\le \gamma_j\ \forall j;\ x_{ij}\in\{0,1\}\big\}$ e $x^*$ è ottima se $\sum_{i,j}x^*_{ij}\ge \sum_{i,j}x_{ij}$ per ogni $x\in X$.

\begin{oss}[cosa può succedere]
\begin{itemize}[nosep]
\item \textbf{Ottimo non unico}: si parla di \emph{una} soluzione ottima, non \emph{la}. Il \emph{valore} ottimo invece è unico. Es.: $\min\{x^4-x^2\}$ ha ottimi $x=\pm\frac{\sqrt2}{2}$, entrambi di valore $-\tfrac14$. Nelle sale operatorie esistono molte assegnazioni diverse con 7 pazienti.
\item \textbf{Problema inammissibile}: $X=\emptyset$. Es.: $\min\{x_1-x_2: x_1+x_2\ge 2,\ 0\le x_1,x_2\le 1\}$.
\item \textbf{Problema illimitato}: per un problema di minimo, per ogni $M$ esiste $x\in X$ con $f(x)< M$ (l'obiettivo scende indefinitamente). Es.: $\max\{2x_1+x_2: x_2\ge x_1,\ x\ge 0\}$.
\item \textbf{Nessun ottimo pur essendo limitato}: $\min\{2x: x>4\}$ ha estremo inferiore 8 non raggiunto. Per questo \textbf{nei modelli non si usano mai disuguaglianze strette}.
\end{itemize}
\end{oss}

\textbf{Min e max sono intercambiabili}:
\[\min\{f(x):x\in X\}=-\max\{-f(x):x\in X\},\]
e le soluzioni ottime dei due problemi coincidono.

\subsection{Forma parametrica, forma compatta (matriciale), istanza}

Uno stesso modello si può scrivere in tre modi:
\begin{itemize}[nosep]
  \item \textbf{forma esplicita (estesa)}: con i numeri dell'istanza, un vincolo per riga (es.\ $3000x_L+5000x_P$\dots);
  \item \textbf{forma parametrica}: con insiemi, indici e parametri simbolici ($\sum_{i\in I}t_ix_{ij}\le\gamma_j\ \forall j\in J$). È quella richiesta all'esame e quella che si scrive nel linguaggio di modellazione del laboratorio: il modello è separato dai dati;
  \item \textbf{forma compatta (matriciale)}: $\max\{c^Tx: Ax\le b,\ x\ge0\}$, dove $c\in\R^n$ è il vettore dei costi/profitti, $A\in\R^{m\times n}$ la matrice dei coefficienti e $b\in\R^m$ il vettore dei termini noti. Serve per la teoria (simplesso, dualità).
\end{itemize}
Per passare alla forma compatta si ``srotolano'' le variabili in un unico vettore. Nel caso del coltivatore (\S\ref{sec:colt}) con $x=(x_L,x_P)^T$:
\[
c=\begin{pmatrix}3000\\5000\end{pmatrix},\quad
A=\begin{pmatrix}1&1\\7&0\\0&3\\10&20\end{pmatrix},\quad
b=\begin{pmatrix}12\\70\\18\\160\end{pmatrix}.
\]
Con variabili a due indici $x_{ij}$ (sale operatorie) il vettore $x$ ha $|I|\cdot|J|$ componenti e ogni vincolo è una riga di $A$.

\section{Tipi di modelli di Programmazione Matematica}

\begin{defbox}[Programmazione Lineare]
Un modello è di \textbf{Programmazione Lineare} se vincoli e funzione obiettivo sono \emph{funzioni lineari} delle variabili: somme di variabili moltiplicate per coefficienti noti. Niente potenze diverse da 1, niente prodotti tra variabili, niente $|\cdot|$, $\max$, $\min$, rapporti tra variabili, ecc.
\end{defbox}

\begin{center}
\begin{tabular}{@{}ll@{}}
\toprule
Variabili & Classe \\ \midrule
tutte continue & Programmazione Lineare (PL, \emph{LP})\\
tutte intere & Programmazione Lineare Intera (PLI, \emph{ILP}); se $\{0,1\}$: binaria\\
alcune intere, altre continue & Programmazione Lineare Mista Intera (PLMI, \emph{MILP})\\
\bottomrule
\end{tabular}
\end{center}

\textbf{Altri tipi} (non trattati nel corso):
\begin{itemize}[nosep]
  \item \textbf{Programmazione non lineare} (eventualmente mista intera): obiettivo o vincoli non lineari. Es.: posizionare magazzini farmaceutici a distanza euclidea $\le K$ km da un gruppo di ospedali: le variabili sono le coordinate $(u,v)$ e il vincolo $\sqrt{(u-a_h)^2+(v-b_h)^2}\le K$ non è lineare.
  \item \textbf{Programmazione stocastica}: alcuni dati sono incerti (es.\ durate degli interventi che sono variabili aleatorie).
  \item \textbf{Simulazione}: emula il comportamento dinamico di un sistema (analisi \emph{what-if}).
  \item \textbf{Teoria dei giochi}: più decisori (es.\ agenti di un mercato).
  \item \textbf{Ottimizzazione multi-criterio}: più obiettivi da valutare contemporaneamente.
\end{itemize}

\begin{oss}[perché insistiamo sulla linearità]
Un problema di PL con milioni di variabili e vincoli si risolve in minuti. Appena si introducono variabili intere (PLI/PLMI) la dimensione trattabile scende a migliaia/decine di migliaia; con vincoli non lineari e non convessi a centinaia. Per questo nella modellazione l'obiettivo è \textbf{scrivere tutto in forma lineare}, introducendo variabili intere/binarie solo quando serve (e servono spesso: decisioni sì/no, logica, costi fissi). Il Cap.~\ref{cap:mod} è in buona parte una raccolta di ``trucchi'' per linearizzare.
\end{oss}

\begin{extra}[rilassamento (appunti a.a.\ precedente)]
Concetto che diventa centrale nella PLI (branch-and-bound) ma conviene averlo chiaro subito.

Dato $P:\min\{f(x):x\in X\}$, un \textbf{rilassamento} è un problema $\bar P:\min\{\bar f(x):x\in \bar X\}$ con
(1) $\bar X\supseteq X$ e (2) $\bar f(x)\le f(x)$ per ogni $x\in X$ (per un massimo: $\bar f(x)\ge f(x)$).

Proprietà (per un minimo):
\begin{itemize}[nosep]
\item il valore ottimo di $\bar P$ è un \textbf{limite inferiore} (\emph{lower bound}) di quello di $P$: $\bar z\le z^*$ (per un massimo è un limite superiore);
\item una soluzione ottima di $\bar P$ può non essere ammissibile per $P$;
\item se una soluzione ottima $\bar x$ di $\bar P$ è ammissibile per $P$ e $\bar f(\bar x)=f(\bar x)$, allora $\bar x$ è ottima anche per $P$.
\end{itemize}
Il rilassamento più importante del corso è il \textbf{rilassamento continuo}: si sostituisce $x\in\{0,1\}$ con $x\in[0,1]$ (o $x\in\Z$ con $x\in\R$). Es.: per lo zaino di \S\ref{sec:zaino} il rilassamento continuo vale 5550, l'ottimo intero 5300.

Una \textbf{restrizione} invece restringe la regione ammissibile ($\hat X\subset X$): fornisce un limite superiore per un minimo. Attenzione: restrizione non è il ``contrario'' di rilassamento.
\end{extra}

\begin{extra}[richiami di convessità (serviranno per la PL)]
\begin{itemize}[nosep]
\item Combinazione di $x',x''$ con pesi $\lambda',\lambda''$: \emph{lineare} (pesi qualsiasi), \emph{affine} ($\lambda'+\lambda''=1$), \emph{conica} ($\lambda',\lambda''\ge0$), \emph{convessa} (entrambe). Lo stesso vale per $m$ vettori.
\item Un insieme $S$ è \textbf{convesso} se contiene il segmento tra due qualsiasi suoi punti: $\lambda x'+(1-\lambda)x''\in S$ per ogni $\lambda\in[0,1]$. L'\textbf{inviluppo convesso} $\mathrm{conv}(V)$ è il più piccolo convesso che contiene $V$.
\item $f$ è \textbf{convessa} se $f(\lambda x'+(1-\lambda)x'')\le \lambda f(x')+(1-\lambda)f(x'')$; concava se vale $\ge$. Le funzioni lineari sono sia convesse sia concave. Il massimo di funzioni convesse è convesso (il minimo no).
\item Se $f$ è convessa, $\{x: f(x)\le\alpha\}$ è convesso. Quindi ogni vincolo lineare definisce un semispazio convesso e la regione ammissibile di una PL, $\{x:Ax\le b,x\ge0\}$, è un \textbf{poliedro} (intersezione di semispazi) convesso.
\item Un insieme $\{(x,y): x\le Cy,\ y\in\{0,1\}\}$ \emph{non} è convesso: le variabili intere distruggono la convessità, ed è questo che rende la PLI difficile.
\end{itemize}
\end{extra}

\chapter{Modelli di Programmazione Lineare}\label{cap:mod}

\section{Come si costruisce un modello}

\begin{center}
\begin{tabular}{@{}lcl@{}}
\toprule
Dati & $\longrightarrow$ & insiemi e parametri\\
Decisioni & $\longrightarrow$ & variabili decisionali\\
Requisiti & $\longrightarrow$ & vincoli\\
Obiettivo & $\longrightarrow$ & funzione obiettivo\\
\bottomrule
\end{tabular}
\end{center}

\begin{esame}
Metodo che conviene seguire \emph{sempre} (e che viene valutato):
\begin{enumerate}[nosep]
\item \textbf{Insiemi e parametri}: elencali con il loro significato e unità di misura.
\item \textbf{Variabili}: per ognuna scrivi \emph{significato}, \emph{indici} (``$\forall i\in I$'') e \emph{dominio} ($\ge0$, libera, intera, $\{0,1\}$). Chiediti: ``se conosco il valore di queste variabili, so ricostruire \emph{tutta} la decisione?''.
\item \textbf{Vincoli}: uno per requisito, con il quantificatore corretto ($\forall$). Controlla che descrivano \emph{tutte e sole} le soluzioni ammissibili: se ne mancano, il modello accetta soluzioni assurde; se sono troppo stretti, si tagliano soluzioni buone (anche l'ottimo).
\item \textbf{Obiettivo}: $\min$ o $\max$ di un'espressione lineare.
\item \textbf{Verifica}: nessun prodotto tra variabili, nessuna disuguaglianza stretta, nessun indice ``libero'' non quantificato, dimensioni coerenti.
\end{enumerate}
Distingui sempre ciò che \emph{decidi} (variabili) da ciò che è \emph{dato}, anche se incerto (parametri: prezzi futuri, meteo, domanda\ldots).
\end{esame}

\section{Il problema del coltivatore}\label{sec:colt}

\textbf{Testo.} Un coltivatore ha 12 ettari da coltivare a lattuga o patate. Risorse: 70 kg di semi di lattuga, 18 t di tuberi, 160 t di stallatico. Resa: 3000 \euro/ha (lattuga), 5000 \euro/ha (patate). Consumi per ettaro: lattuga 7 kg di semi e 10 t di stallatico; patate 3 t di tuberi e 20 t di stallatico. Quanto terreno destinare a ciascuna coltura per massimizzare la resa?


\textbf{Variabili}: $x_L,x_P\ge 0$ = ettari coltivati a lattuga / patate (continue).
\begin{modello}{Coltivatore (forma esplicita)}
\begin{align*}
\max\ & 3000x_L+5000x_P\\
\st\ & x_L+x_P\le 12 &&\text{(terreno)}\\
& 7x_L\le 70 &&\text{(semi di lattuga)}\\
& 3x_P\le 18 &&\text{(tuberi)}\\
& 10x_L+20x_P\le 160 &&\text{(stallatico)}\\
& x_L,x_P\ge 0
\end{align*}
\end{modello}

\textbf{Forma parametrica.} $C$ insieme delle colture, $R$ insieme delle risorse (senza il terreno), $p_j$ resa per ettaro della coltura $j$, $a_{ij}$ quantità di risorsa $i$ usata per ettaro di coltura $j$, $b_i$ disponibilità di $i$, $S$ ettari totali:
\[
\max \sum_{j\in C}p_jx_j\quad \st\ \sum_{j\in C}x_j\le S;\quad \sum_{j\in C}a_{ij}x_j\le b_i\ \ \forall i\in R;\quad x_j\ge0\ \ \forall j\in C.
\]
È il prototipo dei \textbf{problemi di mix produttivo}: si sceglie quanto produrre di ogni bene con risorse limitate.

\paragraph{Soluzione grafica.} Con due variabili si può disegnare la regione ammissibile (un poligono convesso) e ``spingere'' la retta di livello $3000x_L+5000x_P=\alpha$ nella direzione del gradiente $c=(3000,5000)$ finché tocca ancora la regione.

\begin{center}
\begin{tikzpicture}[scale=0.55,>=Stealth]
\fill[blue!12] (0,0)--(10,0)--(10,2)--(8,4)--(4,6)--(0,6)--cycle;
\draw[->] (-0.5,0)--(15,0) node[below] {$x_L$};
\draw[->] (0,-0.5)--(0,10) node[above] {$x_P$};
\foreach \x in {2,4,...,12} \draw (\x,0.1)--(\x,-0.1) node[below,font=\scriptsize] {\x};
\foreach \y in {2,4,6,8} \draw (0.1,\y)--(-0.1,\y) node[left,font=\scriptsize] {\y};
\draw[thick,red!70!black] (12,0)--(2,10) node[above,font=\scriptsize] {terreno};
\draw[thick,green!50!black] (10,-0.3)--(10,9) node[above,font=\scriptsize] {semi};
\draw[thick,orange!80!black] (-0.3,6)--(13,6) node[right,font=\scriptsize] {tuberi};
\draw[thick,violet] (0,8)--(13,1.5) node[above right,font=\scriptsize] {stallatico};
\draw[dashed] (0,8.8) node[left,font=\scriptsize] {$f=44000$}--(14.67,0);
\fill (8,4) circle (4pt) node[above right,font=\small] {$x^*=(8,4)$};
\draw[->,thick] (2,2)--(3.5,4.5) node[above,font=\scriptsize] {$\nabla f$};
\end{tikzpicture}
\end{center}

I vertici della regione sono $(0,0),(10,0),(10,2),(8,4),(4,6),(0,6)$, con valori $0$, $30000$, $40000$, $\mathbf{44000}$, $42000$, $30000$. L'ottimo è $x^*=(8,4)$: 8 ettari a lattuga e 4 a patate, resa 44\,000 \euro. In $x^*$ sono \emph{attivi} (soddisfatti all'uguaglianza) i vincoli di terreno e stallatico; semi e tuberi avanzano (14 kg e 6 t).

\begin{oss}[anticipazione]
L'ottimo di una PL, se esiste, si trova sempre in un \textbf{vertice} della regione ammissibile. Il metodo del simplesso sfrutta questo fatto spostandosi di vertice in vertice adiacente migliorando l'obiettivo. Con variabili intere questo non vale più: il vertice ottimo può avere coordinate frazionarie.
\end{oss}

\section{Il problema della dieta}\label{sec:dieta}

\textbf{Testo.} Mensa scolastica: si vuole il menù giornaliero a costo minimo che fornisca almeno $d=700$ calorie. Dati per etto:

\begin{center}
\begin{tabular}{@{}lcccccc@{}}
\toprule
Cibo & Spaghetti & Riso & Lonza & Cavolo & Insalata & Mela\\ \midrule
Prezzo $p_i$ (\euro/hg) & 0.20 & 0.50 & 2.20 & 0.10 & 0.15 & 0.20\\
Calorie $c_i$ (/hg) & 300 & 250 & 200 & 70 & 180 & 100\\
\bottomrule
\end{tabular}
\end{center}

\begin{modello}{Dieta (forma parametrica)}
$F$: insieme dei cibi; $x_i\ge0$: etti del cibo $i\in F$.
\begin{align*}
\min\ & \sum_{i\in F}p_ix_i\\
\st\ & \sum_{i\in F}c_ix_i\ge d\\
& x_i\ge 0 && \forall i\in F
\end{align*}
\end{modello}
Con più nutrienti $k\in K$ (proteine, grassi, \ldots) con fabbisogno $d_k$ e contenuto $a_{ki}$ si ha un vincolo per nutriente: $\sum_i a_{ki}x_i\ge d_k\ \forall k\in K$. È il prototipo dei \textbf{problemi di copertura della domanda a costo minimo} (duale ``naturale'' del mix produttivo).

\begin{esbox}[soluzione e intuizione]
Con un solo vincolo conviene usare solo il cibo con il miglior rapporto calorie/prezzo: spaghetti $300/0.2=1500$ cal/\euro{} (insalata 1200, cavolo 700, riso e mela 500, lonza 91). Ottimo: $x_S=700/300=2.\overline{3}$ hg, costo $0.4\overline{6}$ \euro. Una dieta di soli spaghetti: il modello fa esattamente ciò che gli chiediamo, non ciò che intendiamo. Servono altri vincoli (qualità, varietà), che introduciamo ora.
\end{esbox}

\subsection{Vincoli di qualità (blending)}\label{sec:blending}

\emph{Almeno il 30\% delle calorie deve provenire da frutta e verdura.} Sia $V=\{\text{cavolo, insalata, mela}\}\subseteq F$. La traduzione naturale
\[ \frac{\sum_{i\in V}c_ix_i}{\sum_{i\in F}c_ix_i}\ge 0.3 \]
\textbf{non è lineare} (rapporto tra variabili). Ma il denominatore è positivo (le calorie totali sono $\ge 700>0$), quindi si può moltiplicare:
\[ \boxed{\sum_{i\in V}c_ix_i\ \ge\ 0.3\sum_{i\in F}c_ix_i} \qquad\Longleftrightarrow\qquad \sum_{i\in V}0.7\,c_ix_i-\sum_{i\in F\setminus V}0.3\,c_ix_i\ge 0. \]
Questo è un \textbf{vincolo di blending} (miscelazione): impone che una componente rappresenti almeno/al più una certa frazione del totale. Compare ovunque ci siano miscele (benzine, leghe, mangimi, portafogli).

Nuovo ottimo: $x_S=1.6\overline{3}$ hg (490 cal) e $x_{\text{insalata}}=1.1\overline{6}$ hg (210 cal $=30\%$), costo $\approx0.502$ \euro.

\begin{oss}
Attenzione al caso in cui il denominatore può annullarsi: se il totale può essere 0, la forma moltiplicata è ancora corretta (diventa $0\ge0$), mentre il rapporto non è definito. La forma lineare è quindi anche più robusta.
\end{oss}

\subsection{Variabili logiche e vincoli di big-\texorpdfstring{$M$}{M}}\label{sec:bigM}

\emph{Nel menù non possono esserci sia gli spaghetti sia il riso.} Una condizione del tipo ``$x_S>0$ e $x_R>0$ non entrambi'' non è esprimibile con le sole variabili continue. Introduciamo \textbf{variabili logiche} (binarie):
\[ y_S\in\{0,1\}: 1 \text{ se gli spaghetti fanno parte della dieta};\qquad y_R\in\{0,1\}: \text{idem per il riso.}\]
Il vincolo logico è semplicemente $y_S+y_R\le 1$. Ora bisogna \textbf{collegare} $x$ e $y$:
\[ x_S>0\Rightarrow y_S=1\quad(\text{equivalente a } y_S=0\Rightarrow x_S=0)\qquad\leadsto\qquad x_S\le M y_S,\quad x_R\le M y_R, \]
dove $M$ è una costante \emph{abbastanza grande da non limitare} i valori sensati di $x$. Se $y_S=0$ il vincolo forza $x_S\le0$, cioè $x_S=0$; se $y_S=1$ diventa $x_S\le M$, inattivo.

\begin{oss}[cosa il big-$M$ non impone]
$x_S\le My_S$ realizza solo l'implicazione $x_S>0\Rightarrow y_S=1$. Non impone $y_S=1\Rightarrow x_S>0$: il modello può porre $y_S=1$ con $x_S=0$. Di solito va bene perché $y$ ha un costo o compare in vincoli restrittivi (quindi all'ottimo non viene accesa inutilmente). Se serve anche l'implicazione inversa si usa un lotto minimo (sotto) con una soglia $\varepsilon>0$ o $q>0$.
\end{oss}

\begin{esame}
\textbf{Come scegliere $M$.} $M$ deve essere un \emph{limite superiore valido} per la variabile (o per l'espressione) che controlla, in ogni soluzione che ci interessa; per il resto conviene che sia il \emph{più piccolo possibile}. Esempi: la capacità di una macchina; il budget diviso per il prezzo; nella dieta, spendere più di quanto costa la dieta di soli spaghetti non ha senso\ldots Un $M$ enorme ($10^9$) è formalmente corretto, ma rende debole il rilassamento continuo (la $y$ può valere $10^{-9}$) e crea problemi numerici ai solutori. All'esame giustifica il valore scelto.
\end{esame}

\subsection{Quantità minima (lotto minimo)}

\emph{Se si mette il riso nel menù, se ne devono mettere almeno $q$ etti.}
\[ q\,y_R\le x_R\le M\,y_R. \]
Se $y_R=0$: $x_R=0$. Se $y_R=1$: $q\le x_R\le M$. L'insieme dei valori ammessi per $x_R$ è $\{0\}\cup[q,M]$, che \textbf{non è convesso}: con sole variabili continue non si potrebbe rappresentare. Questo è il vincolo di \textbf{quantità (lotto) minima}, tipico della produzione (``se accendo la macchina, produco almeno $q$'').

\subsection{Catalogo dei vincoli logici}\label{sec:logici}

Nell'esempio della dieta: al più una fonte di carboidrati $y_S+y_R\le 1$; almeno una fonte di vitamine $y_C+y_I+y_M\ge1$; esattamente una verdura $y_C+y_I=1$.

In generale, con $y_a,y_b,y_c\in\{0,1\}$ associate alle scelte $a,b,c$:

\begin{center}
\renewcommand{\arraystretch}{1.25}
\begin{tabular}{@{}ll@{}}
\toprule
Requisito logico & Vincolo lineare\\ \midrule
al più una tra $a$ e $b$ & $y_a+y_b\le 1$\\
esattamente una tra $a$ e $b$ & $y_a+y_b=1$\\
almeno una tra $a$ e $b$ ($a\lor b$ vera) & $y_a+y_b\ge 1$\\
al più / almeno / esattamente $k$ tra $n$ & $\sum_i y_i\le k$ / $\ge k$ / $=k$\\
non $a$ & $1-y_a$\\
se $a$ allora $b$ ($a\Rightarrow b$) & $y_a\le y_b$\\
se $a$ allora non $b$ & $y_a+y_b\le 1$\\
$a$ se e solo se $b$ & $y_a=y_b$\\
se $a$ e $b$ allora $c$ & $y_a+y_b-1\le y_c$\\
se almeno una tra $a$ e $b$ allora $c$ & $y_a+y_b\le 2y_c$ \quad (oppure $y_a\le y_c,\ y_b\le y_c$)\\
se almeno $m$ tra $y_1,\dots,y_n$ allora $w$ & $\sum_{i=1}^n y_i-m+1\le (n-m+1)\,w$\\
\bottomrule
\end{tabular}
\end{center}

\textbf{Verifica dell'ultimo.} Se $\sum y_i\ge m$, il lato sinistro è $\ge1$ e quindi $w$ deve valere 1. Se $\sum y_i\le m-1$, il lato sinistro è $\le0$ e $w$ è libera. Il lato sinistro vale al più $n-m+1$ (tutte a 1), quindi il coefficiente $n-m+1$ è esattamente quello che serve (è il ``big-$M$'' più piccolo). Con $m=2,n=2$ si ritrova ``se $a$ e $b$ allora $c$''; con $m=1,n=2$ si ritrova ``se almeno una allora $c$''.

\textbf{Metodo generale} per verificare un vincolo logico: fai la tabella di verità, cioè controlla tutte le combinazioni di valori 0/1 e verifica che il vincolo escluda esattamente quelle vietate.

\begin{extra}[variabili che \emph{valgono} un'espressione logica]
A volte serve una variabile $z$ che \emph{sia uguale} a una funzione logica (non solo un vincolo che la imponga vera):
\begin{itemize}[nosep]
\item $z=y_1\land y_2$ (AND): $z\le y_1,\ z\le y_2,\ z\ge y_1+y_2-1$. (Tentazione da evitare: $z=y_1y_2$, prodotto non lineare. Questo è il modo standard di \emph{linearizzare il prodotto di due binarie}.)
\item $z=y_1\lor y_2$ (OR): $z\ge y_1,\ z\ge y_2,\ z\le y_1+y_2$.
\item $z=y_1\oplus y_2$ (XOR): $z\ge y_1-y_2,\ z\ge y_2-y_1,\ z\le y_1+y_2,\ z\le 2-y_1-y_2$.
\end{itemize}
Se $z$ compare solo in un vincolo o nell'obiettivo con un ``verso'' preciso, spesso basta metà dei vincoli (quelli che impediscono a $z$ di ``barare'' nella direzione favorevole).
\end{extra}

\subsection{Attivare e disattivare un vincolo}

Si vuole imporre $\sum_j a_jx_j\le b$ \emph{solo se} $y=1$:
\[ \sum_j a_jx_j\le b+M(1-y). \]
Con $y=1$ è il vincolo originale; con $y=0$ diventa ridondante se $M\ge \max_x\sum_ja_jx_j-b$.

\begin{esbox}[calcolo di $M$]
Vincolo $5x_1+4x_2+7x_3-9x_4\le16$ con $x_1\in\{0,1\}$, $x_2\in[0,5]$, $x_3\in\{0,\dots,4\}$, $x_4\in[0,1]$. Il massimo del lato sinistro si ottiene con le variabili a coefficiente positivo al loro massimo e $x_4$ al minimo: $5+20+28-0=53$. Serve $16+M\ge53$, cioè $M\ge 37$: il valore migliore è $M=37$.
\end{esbox}

\textbf{Vincoli disgiuntivi} (``almeno uno tra due vincoli deve valere''): $g_1(x)\le b_1+M_1(1-y)$, $g_2(x)\le b_2+M_2y$ con $y\in\{0,1\}$. Più in generale, ``almeno $k$ tra $m$ vincoli'': $g_i(x)\le b_i+M_i(1-y_i)$, $\sum_i y_i\ge k$.

\textbf{Costo fisso} (\emph{fixed charge}): costo $f+cx$ se $x>0$, $0$ se $x=0$. Si modella con $y\in\{0,1\}$, obiettivo $fy+cx$ e vincolo $x\le My$. Qui l'implicazione inversa non serve: minimizzando i costi, il solutore non pagherà mai $f$ inutilmente.

\section{Investimento di capitale e problema dello zaino}\label{sec:zaino}

\textbf{Investimento di capitale.} Capitale $B$, $n$ strumenti finanziari; lo strumento $i$ si può acquistare fino a una quantità (in euro) $w_i$ e, se acquistato per intero, stacca una cedola $p_i$. Massimizzare il ritorno totale.

Variabile: $x_i\in[0,1]$ = frazione dello strumento $i$ acquistata (il ritorno è proporzionale: $p_ix_i$).
\begin{modello}{Investimento (continuo)}
\[ \max\sum_{i=1}^n p_ix_i\quad\st\ \sum_{i=1}^n w_ix_i\le B,\qquad 0\le x_i\le 1\ \ \forall i. \]
\end{modello}
(Equivalentemente $z_i\in[0,w_i]$ = euro investiti in $i$, con ritorno $\frac{p_i}{w_i}z_i$.)

\textbf{Variante:} gli strumenti non sono frazionabili, o si compra tutto o niente: $x_i\in[0,1]\ \to\ x_i\in\{0,1\}$. Si ottiene il \textbf{problema dello zaino} (\emph{knapsack}).

\begin{modello}{Problema dello zaino (0/1)}
Dati: oggetti $I$, peso $w_i$ e profitto $p_i$ per $i\in I$, capacità $B$. $x_i=1$ se l'oggetto $i$ è nello zaino.
\[ \max\sum_{i\in I}p_ix_i\quad\st\ \sum_{i\in I}w_ix_i\le B,\qquad x_i\in\{0,1\}\ \ \forall i\in I. \]
\end{modello}

\begin{oss}[continuo vs intero: una differenza enorme]
La versione continua (\emph{zaino frazionario}) si risolve con un algoritmo \emph{greedy}: si ordinano gli oggetti per rapporto $p_i/w_i$ decrescente e si riempie lo zaino in quest'ordine, frazionando solo l'ultimo oggetto che non entra. Complessità $O(n\log n)$. La versione 0/1 invece è \textbf{NP-completa} (Cap.~\ref{cap:compl}): il greedy può sbagliare e non si conosce un algoritmo polinomiale.
\end{oss}

\begin{esbox}[zaino (dagli appunti a.a.\ precedente)]
Capacità 20 kg.
\begin{center}\small
\begin{tabular}{@{}l*{8}{c}@{}}
\toprule
Oggetto & 1&2&3&4&5&6&7&8\\ \midrule
Valore & 1250&300&100&2000&900&300&250&1800\\
Peso & 7.5&1.5&4&8&3&0.5&3&5.5\\
$p/w$ & 167&200&25&250&300&600&83&327\\
\bottomrule
\end{tabular}
\end{center}
\textbf{Rilassamento continuo} (greedy): ordine $6,8,5,4,2,1,\dots$; si prendono $6,8,5,4,2$ (18.5 kg, valore 5300) e $1.5/7.5=0.2$ dell'oggetto 1 (valore 250): \textbf{5550}.
\textbf{Ottimo intero}: $\{2,4,5,6,8\}$, peso 18.5, valore \textbf{5300}. Qui il greedy (senza frazionare) trova l'ottimo, ma in generale non è così: con capacità 10 e oggetti $(p,w)=(6,6),(5,5),(5,5)$ il greedy prende solo il primo (6) mentre l'ottimo è 10.
Il valore del rilassamento (5550) è un \emph{limite superiore}: ci dice subito che nessuna soluzione intera può superare 5550.
\end{esbox}

\section{Obiettivi bottleneck: minimizzare il massimo}\label{sec:bottleneck}

\subsection{Apertura di ambulatori}
Villaggi $I=\{1,\dots,n\}$, siti candidati $J=\{1,\dots,m\}$. Aprire un ambulatorio in $j$ costa $f_j$; budget $B$. Ogni villaggio va assegnato a un ambulatorio aperto; se $i$ è assegnato a $j$ gli abitanti percorrono $p_{ij}$ km. Obiettivo: \emph{minimizzare il percorso più lungo}.

Variabili: $y_j\in\{0,1\}$ (apro in $j$); $x_{ij}\in\{0,1\}$ (assegno $i$ a $j$). L'obiettivo ``naturale'' è
\[ \min\ \max_{i\in I}\ \sum_{j\in J}p_{ij}x_{ij}, \]
non lineare. Si introduce una variabile ausiliaria $z$ (continua) che rappresenta il massimo:

\begin{modello}{Localizzazione con obiettivo bottleneck}
\begin{align*}
\min\ & z\\
\st\ & z\ge \sum_{j\in J}p_{ij}x_{ij} && \forall i\in I &&\text{($z$ maggiora ogni distanza)}\\
& \sum_{j\in J}x_{ij}=1 && \forall i\in I &&\text{(ogni villaggio assegnato)}\\
& x_{ij}\le y_j && \forall i\in I,j\in J &&\text{(solo a siti aperti)}\\
& \sum_{j\in J}f_jy_j\le B &&&&\text{(budget)}\\
& x_{ij},y_j\in\{0,1\}
\end{align*}
\end{modello}
(Alternativa equivalente per il primo vincolo: $z\ge p_{ij}x_{ij}\ \forall i,j$.)

\begin{prop}[linearizzazione del min-max]
Se $f_1,\dots,f_k$ sono lineari, i problemi
\[ \min_{x\in X}\max\{f_1(x),\dots,f_k(x)\}\qquad\text{e}\qquad \min\{z:\ z\ge f_h(x)\ \forall h=1,\dots,k,\ x\in X\} \]
hanno le stesse soluzioni ottime (in $x$) e lo stesso valore ottimo.
\end{prop}
\begin{dimo}
$z\ge\max_h f_h(x)$ equivale a $z\ge f_h(x)$ per ogni $h$. Il secondo problema è quindi $\min\{z: z\ge\max_hf_h(x), x\in X\}$. In una soluzione ottima $(x^*,z^*)$ deve essere $z^*=\max_hf_h(x^*)$: se fosse $z^*>\max_hf_h(x^*)$, la soluzione $(x^*,\max_hf_h(x^*))$ sarebbe ammissibile e strettamente migliore. Quindi all'ottimo $z$ coincide con il massimo e i due problemi si equivalgono. \hfill$\square$
\end{dimo}

\begin{esame}
Il trucco funziona perché si combinano operatori ``opposti'': \textbf{min-max} (e simmetricamente \textbf{max-min}: $\max z$, $z\le f_h(x)\ \forall h$). \textbf{Non} funziona per max-max o min-min: in quel caso $z$ andrebbe all'infinito o servirebbero variabili binarie per scegliere quale $f_h$ è quella ``attiva''. Esempio: $\max\ \max_h f_h(x)$ richiede $z\le f_h(x)+M(1-u_h)$, $\sum_h u_h=1$, $u_h\in\{0,1\}$.
\end{esame}

\textbf{Valore assoluto.} $|f(x)|=\max\{f(x),-f(x)\}$, quindi $\min |f(x)|$ si linearizza con $\min z$, $z\ge f(x)$, $z\ge -f(x)$. Esempio (appunti): suddividere oggetti di valore $v_i$ in due gruppi il più equilibrati possibile, $x_i=1$ se $i$ va nel gruppo 1:
\[ \min z\quad \st\ z\ge 2\sum_iv_ix_i-\sum_iv_i,\quad z\ge -2\sum_iv_ix_i+\sum_iv_i,\quad x_i\in\{0,1\}. \]

\section{Variabili libere}\label{sec:libere}

Non tutte le variabili sono non negative. \textbf{Esempio (orario dei treni).} Treni $T=\{1,\dots,n\}$, stazioni $S=\{1,\dots,m\}$, orario di arrivo $a_{ts}$. Vogliamo studiare l'effetto di una \emph{variazione} degli orari: la variabile naturale è
\[ \delta_{ts}\in\R \quad(\text{libera}):\ \text{variazione dell'orario di arrivo di } t \text{ in } s,\qquad \text{nuovo orario}=a_{ts}+\delta_{ts}, \]
che può essere un anticipo ($\delta<0$) o un ritardo ($\delta>0$). Nei modelli si scrive esplicitamente ``$\delta_{ts}$ libera'' (o $\delta_{ts}\in\R$).

\textbf{Riduzione a variabili non negative.} Ogni variabile libera si scrive come differenza di due variabili non negative:
\[ x=x^+-x^-,\qquad x^+,x^-\ge0. \]
Serve perché il metodo del simplesso lavora con variabili $\ge0$.

\begin{extra}[forma standard di una PL]
Il simplesso richiede la \textbf{forma standard}: $\min c^Tx$, $Ax=b$, $x\ge0$. Ci si arriva sempre:
\begin{itemize}[nosep]
\item $\max f\ \to\ -\min(-f)$;
\item $a^Tx\le b\ \to\ a^Tx+s=b$, $s\ge0$ (variabile di \emph{slack});
\item $a^Tx\ge b\ \to\ a^Tx-s=b$, $s\ge0$ (variabile di \emph{surplus});
\item $x$ libera $\to x^+-x^-$; $x\le0\to x=-x'$ con $x'\ge0$.
\end{itemize}
Esempio: $\max\{3x_1+5x_2-x_3:\ x_2+2x_3\le5,\ 2x_1-x_3=0,\ x_1,x_2\ge0,\ x_3\text{ libera}\}$ diventa
$-\min\{-3x_1-5x_2+x_3^+-x_3^-:\ x_2+2x_3^+-2x_3^-+s_1=5,\ 2x_1-x_3^++x_3^-=0,\ x,s\ge0\}$.
\end{extra}

\begin{oss}[minimizzare scostamenti]
Per minimizzare la variazione totale $\sum|\delta_{ts}|$: $\delta_{ts}=\delta_{ts}^+-\delta_{ts}^-$ e obiettivo $\min\sum(\delta_{ts}^++\delta_{ts}^-)$. All'ottimo al più una tra $\delta^+$ e $\delta^-$ è positiva (altrimenti si potrebbero ridurre entrambe della stessa quantità lasciando $\delta$ invariato e diminuendo l'obiettivo), quindi $\delta^++\delta^-=|\delta|$. Funziona solo se si \emph{minimizza} la somma; massimizzando non va bene.
\end{oss}

\section{Variabili discrete}\label{sec:discrete}

\textbf{Esempio (dimensionamento di una rete di telecomunicazioni).} La capacità di un collegamento si ottiene installando canali di capacità fissata, quindi può assumere solo valori in un insieme finito $\{d_1,\dots,d_r\}$ (non tutti i reali, e nemmeno tutti gli interi).

Per ogni collegamento $e$ si introducono binarie $y_{ek}\in\{0,1\}$, $k=1,\dots,r$ (1 se si sceglie il valore $d_k$):
\[ \sum_{k=1}^r y_{ek}=1,\qquad c_e=\sum_{k=1}^r d_k\,y_{ek}. \]
La variabile $c_e$ (che può anche essere sostituita dall'espressione ovunque compaia) assume esattamente uno dei valori ammessi. Se i canali hanno tutti capacità $u$ e se ne possono installare quanti se ne vuole, basta invece una variabile intera: $c_e=u\,n_e$, $n_e\in\Z_{\ge0}$.

\begin{esame}
Riassunto dei domini: \emph{continue} $x\ge0$ (quantità divisibili); \emph{libere} $x\in\R$ (variazioni, scostamenti, saldi); \emph{intere} $x\in\Z_{\ge0}$ (numero di oggetti indivisibili: infermieri, camion, confezioni); \emph{binarie} $x\in\{0,1\}$ (decisioni sì/no, assegnamenti, logica); \emph{discrete} in $\{d_1,\dots,d_r\}$ (tramite binarie di selezione).
\end{esame}

\section{Altri schemi ricorrenti}

\subsection{Pianificazione multi-periodo con magazzino (bilanciamento)}
Produzione nei periodi $t=1,\dots,T$ con domanda $d_t$, produzione $x_t\le c$, livello di magazzino a fine periodo $s_t$ (con $s_0$ dato), capacità del magazzino $f$ e costo di stoccaggio $g$:
\begin{align*}
\min\ & g\sum_t s_t\\
\st\ & s_{t-1}+x_t=d_t+s_t && \forall t \quad(\text{ciò che entra} = \text{ciò che esce})\\
& 0\le x_t\le c,\quad 0\le s_t\le f && \forall t.
\end{align*}
Il vincolo di bilanciamento (\emph{flow conservation}) è lo schema chiave di tutti i problemi con magazzino (Esercitazione, es.\ 8, 12, 16). Aggiungendo $y_t\in\{0,1\}$ con costo fisso e $qy_t\le x_t\le Qy_t$ si ottiene il \emph{lot sizing}.

\subsection{Copertura, partizione, impaccamento di insiemi}
Dati elementi $i\in M$ e sottoinsiemi $j\in N$ con $a_{ij}=1$ se $i\in S_j$, e $x_j\in\{0,1\}$ = ``scelgo $S_j$'':
\begin{itemize}[nosep]
\item \textbf{copertura} (\emph{set covering}): ogni elemento coperto almeno una volta, $\sum_j a_{ij}x_j\ge1\ \forall i$, $\min\sum_jc_jx_j$ (antenne, ambulanze, uffici);
\item \textbf{partizione} (\emph{set partitioning}): esattamente una volta, $\sum_j a_{ij}x_j=1$;
\item \textbf{impaccamento} (\emph{set packing}): al più una volta, $\sum_ja_{ij}x_j\le1$, $\max\sum_jc_jx_j$.
\end{itemize}
Casi su grafo: \emph{vertex cover} ($x_i+x_j\ge1$ per ogni lato) e \emph{insieme indipendente} ($x_i+x_j\le1$ per ogni lato).

\subsection{Assegnamento}
Ogni elemento $i$ assegnato a esattamente una risorsa $j$: $\sum_jx_{ij}=1\ \forall i$, con eventuali capacità $\sum_iw_ix_{ij}\le W_j\ \forall j$ e collegamento con l'apertura della risorsa $x_{ij}\le y_j$ (oppure la versione aggregata $\sum_iw_ix_{ij}\le W_jy_j$).

\chapter{Cenni di complessità computazionale}\label{cap:compl}

\section{Problemi facili e difficili}

Prima di affrontare un problema di ottimizzazione occorre valutarne la \textbf{difficoltà}, che è legata agli algoritmi in grado di risolverlo:
\begin{itemize}[nosep]
\item \textbf{problema facile}: esiste un algoritmo \emph{polinomiale} che lo risolve;
\item \textbf{problema difficile}: non si conosce un algoritmo polinomiale che lo risolva.
\end{itemize}
La \textbf{complessità computazionale} studia la difficoltà dei problemi; servono due ingredienti: la \emph{dimensione dell'istanza} e la \emph{complessità degli algoritmi}.

\section{Istanza e dimensione dell'istanza}

\begin{defbox}[problema e istanza]
Un \textbf{problema} è una domanda generale, formulata su parametri simbolici. Un'\textbf{istanza} è un esempio specifico del problema, ottenuto fissando i valori numerici dei parametri.
\end{defbox}
Il problema dello zaino è ``dati $I,p,w,B$, risolvi $\max\{\sum p_ix_i:\sum w_ix_i\le B, x\in\{0,1\}^{|I|}\}$''; un'istanza è, per esempio, quella con gli 8 oggetti e $B=20$ di \S\ref{sec:zaino}.

\begin{defbox}[dimensione dell'istanza]
La \textbf{dimensione} di un'istanza è il numero di bit necessari a codificarla (in modo ``ragionevole'', cioè con i numeri in base 2 e non in unario).
\end{defbox}
Un intero $a$ richiede circa $\lceil\log_2(a+1)\rceil$ bit. Per lo zaino con $N=|I|$ oggetti:
\begin{center}
\begin{tabular}{@{}ll@{}}
\toprule
Dato & Bit \\ \midrule
numero di oggetti $N$ & $\lceil\log_2 N\rceil$\\
vettore dei profitti $p$ & $\sum_{i\in I}\lceil\log_2 p_i\rceil\le N\lceil\log_2 p_{\max}\rceil$\\
vettore dei pesi $w$ & $\sum_{i\in I}\lceil\log_2 w_i\rceil\le N\lceil\log_2 w_{\max}\rceil$\\
capacità $B$ & $\lceil\log_2 B\rceil$\\ \midrule
totale & $O\big(N(\log p_{\max}+\log w_{\max})+\log B\big)$\\
\bottomrule
\end{tabular}
\end{center}
Il punto chiave: la dimensione cresce \emph{linearmente} con il numero di oggetti ma solo \emph{logaritmicamente} con il valore dei numeri.

\section{Complessità degli algoritmi}

Un \textbf{algoritmo} è una sequenza finita di operazioni elementari che risolve un problema; per un problema di ottimizzazione, lo risolve se trova una soluzione ottima \emph{per ogni istanza}.

\begin{defbox}[complessità di un algoritmo]
La complessità di un algoritmo è la funzione che, per ogni dimensione $n$ dell'istanza, dà il numero di operazioni elementari (assegnamenti, confronti, operazioni aritmetiche) eseguite \textbf{nel caso peggiore} tra tutte le istanze di dimensione $n$.
\end{defbox}

\begin{esbox}[massima distanza tra $n$ punti]
Data la matrice delle distanze $\{d_{ij}\}$ tra $n$ punti, calcolare la distanza massima.
\begin{center}
\begin{tabular}{@{}ll@{}}
\texttt{dmax := 0} & 1 assegnamento\\
\texttt{per i = 1..n} & $n$ iterazioni\\
\quad\texttt{per j = 1..n} & $n$ iterazioni\\
\qquad\texttt{se d[i][j] > dmax} & 1 confronto\\
\qquad\quad\texttt{dmax := d[i][j]} & (al più) 1 assegnamento
\end{tabular}
\end{center}
Caso peggiore (le distanze arrivano in ordine crescente, quindi si aggiorna sempre): $1+2n^2$ operazioni, cioè $O(n^2)$. Sfruttando la simmetria ($j>i$) si fanno $n(n-1)/2$ iterazioni: ancora $O(n^2)$, cambia solo la costante. Nota che l'input ha $n^2$ numeri, quindi l'algoritmo è lineare nella dimensione dell'input.
\end{esbox}

\begin{defbox}[notazione $O$]
Date $f(n),g(n)$, si dice che $f(n)=O(g(n))$ se esistono $c>0$ e $n_0$ tali che $f(n)\le c\,g(n)$ per ogni $n>n_0$.
\end{defbox}
Si ignorano costanti moltiplicative e termini di ordine inferiore: $3n^2+10n+7=O(n^2)$.
\begin{itemize}[nosep]
\item \textbf{algoritmo polinomiale}: complessità $O(n^d)$ per qualche costante $d$;
\item \textbf{algoritmo esponenziale}: complessità del tipo $O(2^n)$ (o comunque non limitata da alcun polinomio).
\end{itemize}

\begin{center}
\begin{tabular}{@{}lrrrr@{}}
\toprule
$n$ & 10 & 20 & 50 & 100\\ \midrule
$n^2$ & 100 & 400 & 2\,500 & 10\,000\\
$n^3$ & 1\,000 & 8\,000 & 125\,000 & $10^6$\\
$2^n$ & 1\,024 & $\approx10^6$ & $\approx 10^{15}$ & $\approx10^{30}$\\
\bottomrule
\end{tabular}
\end{center}
A $10^9$ operazioni al secondo, $2^{100}$ operazioni richiedono circa $4\cdot10^{13}$ anni. Enumerare tutte le $2^N$ soluzioni dello zaino non è un'opzione.

\begin{extra}[algoritmi pseudo-polinomiali]
Lo zaino con dati interi si risolve con la programmazione dinamica: $F(k,b)$ = miglior profitto usando i primi $k$ oggetti con capacità $b$,
\[F(k,b)=\max\{F(k-1,b),\ F(k-1,b-w_k)+p_k\ (\text{se } w_k\le b)\},\qquad F(0,b)=0.\]
La tabella ha $N\cdot B$ celle: complessità $O(NB)$. Sembra polinomiale ma \textbf{non lo è}: $B$ è codificato con $\log_2B$ bit, quindi $O(NB)=O(N\,2^{\log_2B})$ è esponenziale nella dimensione dell'input. Un algoritmo polinomiale nel \emph{valore} dei numeri (non nella loro lunghezza) si dice \emph{pseudo-polinomiale}. Utile in pratica se $B$ è piccolo; non contraddice la NP-completezza dello zaino.
\end{extra}

\section{Problemi di riconoscimento}

La teoria si sviluppa per problemi con risposta SÌ/NO.
\begin{defbox}[versione di riconoscimento]
La \textbf{versione di riconoscimento} (o decisionale) di un problema di ottimizzazione $\max\{cx: Ax\le b,\ x\ge0,\ x\in\Z^n\}$, data una soglia $k$, chiede:
\[ \exists\,x:\ Ax\le b,\ x\ge 0,\ x\in\Z^n,\ cx\ge k\ ? \]
(per un minimo: $cx\le k$).
\end{defbox}
\textbf{Zaino (riconoscimento):} dati oggetti con profitti $p_i$, pesi $w_i$, capacità $B$ e soglia $k$, esiste un sottoinsieme di oggetti di peso complessivo $\le B$ e profitto complessivo $\ge k$?

\begin{oss}[perché basta studiare il riconoscimento]
Se so risolvere l'ottimizzazione, so rispondere al riconoscimento (confronto il valore ottimo con $k$). Viceversa, con dati interi il valore ottimo sta in un intervallo $[L,U]$ con $\log(U-L)$ polinomiale nella dimensione: con una \emph{ricerca binaria} su $k$ bastano polinomialmente molte chiamate al riconoscimento. Quindi le due versioni sono ``ugualmente difficili'' a meno di un fattore polinomiale.
\end{oss}

\section{Le classi \texorpdfstring{$\Pc$}{P} e \texorpdfstring{$\NP$}{NP}}

\begin{defbox}[classe $\Pc$]
Un problema appartiene a $\Pc$ se esiste un algoritmo che, per ogni istanza del problema di riconoscimento, decide in tempo polinomiale se la risposta è SÌ o NO. (Il corrispondente problema di ottimizzazione si risolve all'ottimo in tempo polinomiale.)
\end{defbox}

\begin{defbox}[classe $\NP$]
Un problema appartiene a $\NP$ se, per ogni istanza con risposta SÌ, esiste un \textbf{certificato} $\gamma(I)$ (tipicamente una soluzione) di dimensione polinomiale e un algoritmo $A$ che verifica \textbf{in tempo polinomiale} che $\gamma(I)$ attesta la risposta SÌ.
\end{defbox}
In parole: \emph{trovare} la soluzione può essere difficile, ma \emph{controllarla} è facile.

\begin{itemize}[nosep]
\item $\NP$ significa \textbf{Non-deterministico Polinomiale} (una macchina che ``indovina'' il certificato e poi lo verifica), \textbf{non} ``non polinomiale''!
\item $\Pc\subseteq\NP$: se il problema è in $\Pc$, il verificatore può ignorare il certificato e risolvere l'istanza da zero in tempo polinomiale.
\item Se $\Pc\subsetneq\NP$ oppure $\Pc=\NP$ è uno dei grandi problemi aperti della matematica (uno dei Millennium Prize Problems). Si congettura $\Pc\ne\NP$.
\end{itemize}

\begin{esbox}[lo zaino è in $\NP$]
Certificato: il sottoinsieme $S$ di oggetti (un vettore $x\in\{0,1\}^N$, dimensione $N$). Verifica: calcolo $\sum_{i\in S}w_i$ e lo confronto con $B$; calcolo $\sum_{i\in S}p_i$ e lo confronto con $k$. Sono $O(N)$ somme e 2 confronti: tempo polinomiale.
\end{esbox}
Notare l'asimmetria: per certificare una risposta NO (``nessun sottoinsieme raggiunge $k$'') non si conosce un certificato breve.

\section{Riduzioni e \texorpdfstring{$\NP$}{NP}-completezza}

\begin{defbox}[riduzione polinomiale]
Dati due problemi $P_1$ e $P_2$, $P_1$ \textbf{si riduce in tempo polinomiale} a $P_2$ ($P_1\red P_2$) se per ogni istanza $I_1$ di $P_1$ si può costruire in tempo polinomiale un'istanza $I_2$ di $P_2$ tale che dalla soluzione di $I_2$ si ricava in tempo polinomiale la soluzione di $I_1$.
\end{defbox}

\begin{center}
\begin{tikzpicture}[>=Stealth, box/.style={draw,rounded corners,minimum height=0.8cm,font=\small,fill=gray!8}]
\node[box] (i1) {$I_1$ (istanza di $P_1$)};
\node[box,right=1.2cm of i1] (i2) {$I_2$};
\node[box,right=1.2cm of i2] (a2) {algoritmo $A_2$};
\node[box,right=1.2cm of a2] (s1) {soluzione di $I_1$};
\draw[->] (i1)--node[above,font=\scriptsize]{poly} (i2);
\draw[->] (i2)--(a2);
\draw[->] (a2)--node[above,font=\scriptsize]{poly} (s1);
\draw[dashed,rounded corners] ($(i1.north west)+(-0.2,0.5)$) rectangle ($(s1.south east)+(0.2,-0.3)$);
\node[font=\scriptsize,fill=white] at ($(i2)!0.5!(a2)+(0,0.8)$) {algoritmo $A_1$ per $P_1$};
\end{tikzpicture}
\end{center}

Conseguenze:
\begin{itemize}[nosep]
\item se esiste un algoritmo $A_2$ per $P_2$, ottengo un algoritmo $A_1$ per $P_1$ che usa $A_2$; se $A_2$ è polinomiale, lo è anche $A_1$;
\item se $P_1\red P_2$, allora \textbf{$P_2$ è almeno tanto difficile quanto $P_1$};
\item se $P_1\red P_2$ e $P_2\in\Pc$, allora $P_1\in\Pc$;
\item (contronominale) se $P_1\red P_2$ e $P_1$ è difficile, allora anche $P_2$ è difficile.
\end{itemize}

\begin{esame}
\textbf{Il verso della riduzione} è l'errore più comune. Per dimostrare che un problema \emph{nuovo} $P$ è difficile si riduce un problema \emph{noto} difficile \emph{a} $P$: $P_{\text{noto}}\red P$ (``se sapessi risolvere $P$ saprei risolvere anche il problema noto''). Ridurre $P$ a un problema difficile ($P\red P_{\text{noto}}$) non dice nulla sulla difficoltà di $P$: tutto si riduce a un problema difficile.
\end{esame}

\begin{defbox}[$\NP$-completo]
Un problema $P$ è \textbf{$\NP$-completo} se
\begin{enumerate}[label=\roman*),nosep]
\item $P\in\NP$;
\item $P_1\red P$ per ogni $P_1\in\NP$.
\end{enumerate}
Cioè $P$ è almeno tanto difficile quanto ogni altro problema in $\NP$. Se vale solo ii), $P$ si dice $\NP$-\emph{difficile} (\emph{NP-hard}); si usa questo termine per le versioni di ottimizzazione, che non sono problemi SÌ/NO.
\end{defbox}
Per i problemi $\NP$-completi non si conosce un algoritmo polinomiale che garantisca l'ottimo per ogni istanza. Se se ne trovasse uno per \emph{un} solo problema $\NP$-completo, allora $\Pc=\NP$.

\textbf{Come si dimostra che $P$ è $\NP$-completo} (senza ridurre ogni problema di $\NP$ uno per uno):
\begin{enumerate}[nosep]
\item si dimostra che $P\in\NP$ (esibendo certificato e verifica polinomiale);
\item si dimostra che $P_1\red P$ per \emph{un} problema $P_1$ già noto come $\NP$-completo.
\end{enumerate}
\begin{prop}
La riducibilità è transitiva: $P'\red P_1$ e $P_1\red P$ implicano $P'\red P$. Quindi i passi 1--2 bastano.
\end{prop}
\begin{dimo}
Transitività: si compongono le due trasformazioni; la composizione di funzioni polinomiali è polinomiale (un polinomio di un polinomio è un polinomio, e l'output di una trasformazione polinomiale ha dimensione polinomiale). Poiché $P_1$ è $\NP$-completo, $P'\red P_1$ per ogni $P'\in\NP$; per transitività $P'\red P$ per ogni $P'\in\NP$; con $P\in\NP$ si conclude. \hfill$\square$
\end{dimo}

\begin{center}
\begin{tikzpicture}[scale=0.9]
\draw[thick,fill=blue!5] (0,0) ellipse (4.8 and 2.2);
\node at (0,1.7) {$\NP$};
\draw[thick,fill=green!12] (-2.1,0) ellipse (1.9 and 1.15);
\node[align=center,font=\small] at (-2.1,0) {$\Pc$\\ PL continua,\\ cammino minimo};
\draw[thick,fill=red!12] (2.2,0) ellipse (1.9 and 1.15);
\node[align=center,font=\small] at (2.2,0) {$\NP$-completi\\ SAT, PLI,\\ zaino};
\node[font=\scriptsize,align=center] at (0,-2.5) {Quadro nell'ipotesi (creduta vera) $\Pc\neq\NP$};
\end{tikzpicture}
\end{center}

\section{Il problema della soddisfacibilità (SAT)}

Il primo problema dimostrato $\NP$-completo (Cook, 1971, con una dimostrazione ``diretta'' che simula una macchina di Turing) è la \textbf{soddisfacibilità}.
\begin{defbox}[SAT]
Date $n$ variabili booleane $y_1,\dots,y_n$ e $m$ espressioni booleane (\emph{clausole}) $C_1,\dots,C_m$, ciascuna disgiunzione ($\lor$) di letterali (variabili o loro negazioni $\bar y_i$), ad esempio $C_1=y_1\lor y_2\lor\bar y_3$, esiste un assegnamento di valori di verità alle variabili che renda \emph{tutte} le clausole vere?
\end{defbox}
Esempio: $(y_1\lor y_2\lor\bar y_3)\land(\bar y_1\lor y_3)\land(\bar y_2\lor\bar y_3)$ è soddisfatta da $y_1=1,y_2=0,y_3=1$.

\section{La Programmazione Lineare Intera è \texorpdfstring{$\NP$}{NP}-completa}

\begin{prop}
La versione di riconoscimento di un generico problema di PLI
\[ \exists x:\ Ax\le b,\ x\ge0,\ x\in\Z^n,\ cx\le k\ ? \]
è $\NP$-completa.
\end{prop}

\textbf{1) PLI $\in\NP$.} Data una soluzione $x$ come certificato, bisogna: controllare tutti i vincoli ($m$ prodotti scalari, $O(mn)$ operazioni) e l'interezza; calcolare $cx$ e confrontarlo con $k$. Tutto in tempo polinomiale. (Sottigliezza tecnica, non richiesta: si dimostra che se esiste una soluzione intera ammissibile ne esiste una con componenti di dimensione polinomiale, quindi il certificato è ``corto''.)

\textbf{2) SAT $\red$ PLI.} Data un'istanza di SAT costruiamo un'istanza di PLI:
\begin{center}
\begin{tabular}{@{}ll@{}}
\toprule
Istanza SAT & Istanza PLI\\ \midrule
variabile booleana $y_i$ & variabile $x_i\in\{0,1\}$ (cioè $0\le x_i\le1$, $x_i\in\Z$)\\
letterale $y_i$ & $x_i$\\
letterale $\bar y_i$ & $1-x_i$\\
clausola $y_1\lor y_2\lor\bar y_3$ vera & $x_1+x_2+(1-x_3)\ge1$\\
\bottomrule
\end{tabular}
\end{center}
In generale, per la clausola $C_j$ con letterali positivi $P_j$ e negati $N_j$:
\[ \sum_{i\in P_j}x_i+\sum_{i\in N_j}(1-x_i)\ge 1\qquad \forall j=1,\dots,m, \]
cioè ``almeno un letterale della clausola è vero''. Obiettivo irrilevante: $c=0$, $k=0$ (la domanda diventa ``esiste una soluzione ammissibile?'').

\emph{Correttezza}: un assegnamento di verità soddisfa tutte le clausole se e solo se il corrispondente vettore 0/1 soddisfa tutti i vincoli. \emph{Tempo}: una variabile per variabile booleana e un vincolo per clausola: la costruzione è lineare nella dimensione dell'istanza SAT; la risposta si trasferisce identica.

\begin{esbox}[istanza costruita]
$(y_1\lor y_2\lor\bar y_3)\land(\bar y_1\lor y_3)\land(\bar y_2\lor\bar y_3)$ diventa
\[ x_1+x_2-x_3\ge0,\qquad -x_1+x_3\ge0,\qquad -x_2-x_3\ge-1,\qquad x\in\{0,1\}^3. \]
$x=(1,0,1)$ è ammissibile, coerentemente con l'assegnamento trovato prima.
\end{esbox}

\begin{esbox}[un'altra riduzione: Subset Sum $\red$ Zaino]
\textsc{Subset Sum} ($\NP$-completo): dati interi $a_1,\dots,a_N$ e $T$, esiste un sottoinsieme con somma esattamente $T$? Costruisco l'istanza di zaino con $p_i=w_i=a_i$, $B=k=T$. Un sottoinsieme ha peso $\le T$ e profitto $\ge T$ se e solo se ha somma esattamente $T$. Costruzione lineare, quindi lo zaino (riconoscimento) è $\NP$-completo.
\end{esbox}

\section{Conseguenze pratiche}

\begin{itemize}
\item \textbf{PL con variabili continue è in $\Pc$} (metodo dell'ellissoide, Khachiyan 1979; metodi a punto interno, Karmarkar 1984). Il simplesso, curiosamente, ha complessità esponenziale nel caso peggiore (esempi di Klee--Minty) ma è velocissimo in pratica.
\item \textbf{PLI e PLMI sono $\NP$-difficili}. Aggiungere il vincolo di interezza cambia radicalmente la difficoltà: è il motivo per cui nei modelli si usano variabili intere solo quando servono.
\item $\NP$-completo non vuol dire ``irrisolvibile'': molte istanze reali si risolvono bene con metodi esatti (branch-and-bound, piani di taglio) che sfruttano i rilassamenti, oppure si usano euristiche (soluzioni buone senza garanzia di ottimalità).
\item Un problema specifico può essere facile anche se la sua formulazione naturale è una PLI (es.\ cammino minimo, assegnamento, flussi su grafo): la PLI è $\NP$-completa come \emph{classe}, non ogni sua istanza o sottoclasse. Si vedrà nella parte sui grafi.
\end{itemize}

\begin{esame}
Domande teoriche tipiche: definire istanza, dimensione, complessità di un algoritmo (caso peggiore); definire $\Pc$, $\NP$, riduzione, $\NP$-completo; spiegare perché $\Pc\subseteq\NP$; mostrare che un problema è in $\NP$ (certificato + verifica); scrivere la riduzione SAT $\red$ PLI; spiegare il verso di una riduzione.
\end{esame}

\chapter[Programmazione Lineare: geometria e ottimo]{Programmazione Lineare: forme, geometria e proprietà dell'ottimo}\label{cap:pl}

\begin{oss}[nota su questo capitolo]
Scritto a partire dalle slide \texttt{04\_ProgrammazioneLineare} (pubblicate il 3/10, prima della lezione). Molte slide sono riquadri vuoti completati alla lavagna: definizioni, esempi e dimostrazioni sono stati ricostruiti in modo standard. Quando avrai visto la lezione, segnala eventuali differenze di notazione e le allineo.
\end{oss}

Parte del programma che inizia qui: risoluzione grafica, forme alternative, geometria della PL, proprietà della soluzione ottima; poi (prossime slide) metodo del simplesso e dualità. In tutto il capitolo le variabili sono \textbf{continue}.

\section{Il coltivatore, riscalato}

Esprimendo la resa in migliaia di euro e lo stallatico in decine di tonnellate, il modello di \S\ref{sec:colt} diventa
\[
\max\ 3x_L+5x_P\quad \st\ x_L+x_P\le12,\ \ x_L\le10,\ \ x_P\le6,\ \ x_L+2x_P\le16,\ \ x_L,x_P\ge0,
\]
cioè $\max\{c^Tx: Ax\le b,\ x\ge0\}$ con
$c=\binom{3}{5}$ e
\[
A=\begin{pmatrix}1&1\\1&0\\0&1\\1&2\end{pmatrix},\qquad b=\begin{pmatrix}12\\10\\6\\16\end{pmatrix}.
\]
(Dividere un vincolo per una costante positiva o riscalare l'obiettivo non cambia le soluzioni ottime; l'ottimo è ancora $(8,4)$, con valore $44$ migliaia di euro.)

\section{Definizioni (richiamo)}
Per $(P)\ \max\{c^Tx: Ax\le b, x\ge0\}$:
\begin{itemize}[nosep]
\item \textbf{soluzione ammissibile}: $\bar x$ con $A\bar x\le b$, $\bar x\ge0$;
\item \textbf{regione ammissibile}: $X=\{x\in\R^n: Ax\le b,\ x\ge0\}$;
\item \textbf{problema inammissibile}: $X=\emptyset$;
\item \textbf{soluzione ottima}: $x^*\in X$ con $c^Tx^*\ge c^Tx$ per ogni $x\in X$;
\item \textbf{problema illimitato}: per ogni $M$ esiste $x\in X$ con $c^Tx>M$.
\end{itemize}

\section{Forme di un problema di PL}

In generale un problema di PL può avere vincoli di tutti i tipi e variabili di tutti i segni:
\[
\max\ (\text{o}\min)\ c^Tx\quad\st\ A_1x\le b_1,\quad A_2x\ge b_2,\quad A_3x=b_3,\quad x_j\ge0,\ x_j\le0\ \text{o libere}.
\]
Due forme ``di riferimento'':
\begin{center}
\begin{tabular}{@{}c@{\qquad\qquad}c@{}}
\textbf{Forma standard} & \textbf{Forma canonica (o generale)}\\[3pt]
$\begin{aligned}\min\ &c^Tx\\ \st\ &Ax=b\\ &x\ge0\end{aligned}$ \quad (con $b\ge0$)
&
$\begin{aligned}\max\ &c^Tx\\ \st\ &Ax\le b\\ &x\ge0\end{aligned}$
\end{tabular}
\end{center}
La forma standard serve al \textbf{simplesso} (lavora con sistemi di equazioni), la forma canonica è comoda per la \textbf{geometria} e per la \textbf{dualità}.

\begin{prop}
Ogni problema di PL può essere scritto in modo equivalente in forma standard, e in modo equivalente in forma canonica.
\end{prop}
``Equivalente'' significa: da ogni soluzione (ottima) dell'uno si ricava in modo immediato una soluzione (ottima) dell'altro, con lo stesso valore dell'obiettivo (a meno del segno).

\subsection{Trasformazioni}
\begin{center}
\renewcommand{\arraystretch}{1.3}
\begin{tabular}{@{}p{0.2\textwidth}p{0.36\textwidth}p{0.36\textwidth}@{}}
\toprule
 & Da & A\\ \midrule
\textbf{Obiettivo} & $\max c^Tx$ & $-\min(-c^T x)$\\
 & $\min c^Tx$ & $-\max(-c^Tx)$\\ \midrule
\textbf{Variabili} & $x_j\le0$ & $x_j=-x_j'$, \ $x_j'\ge0$\\
 & $x_j$ libera & $x_j=x_j^+-x_j^-$, \ $x_j^+,x_j^-\ge0$\\ \midrule
\textbf{Vincoli} & $a^Tx\le b$ & $a^Tx+s=b$, \ $s\ge0$ \ (\emph{slack})\\
 & $a^Tx\ge b$ & $a^Tx-s=b$, \ $s\ge0$ \ (\emph{surplus})\\
 & $a^Tx\ge b$ & $-a^Tx\le -b$\\
 & $a^Tx=b$ & $a^Tx\le b$ \ e \ $-a^Tx\le -b$\\
 & $a^Tx=b$ con $b<0$ & $-a^Tx=-b$ \ (per avere $b\ge0$)\\
\bottomrule
\end{tabular}
\end{center}
Le variabili di slack/surplus hanno un significato: la slack del vincolo di terreno è l'\emph{ettaraggio non usato}. All'ottimo, slack $=0$ significa vincolo \emph{attivo} (risorsa satura).

\begin{esbox}[coltivatore in forma standard]
\[
\begin{aligned}
-\min\ & -3x_L-5x_P\\
\st\ & x_L+x_P+s_1=12\\
& x_L+s_2=10\\
& x_P+s_3=6\\
& x_L+2x_P+s_4=16\\
& x_L,x_P,s_1,s_2,s_3,s_4\ge0
\end{aligned}
\]
All'ottimo $(x_L,x_P)=(8,4)$: $s=(0,2,2,0)$ (terreno e stallatico saturi; avanzano 2 unità di semi e di tuberi nelle unità riscalate). Il coltivatore è \emph{già} in forma canonica.
\end{esbox}

\begin{esbox}[problema misto in entrambe le forme]
$\min\{x_1-2x_2:\ x_1+x_2\ge 2,\ \ x_1-x_2=-1,\ \ x_1\ge0,\ x_2\text{ libera}\}$.
\textbf{Standard}: $x_2=x_2^+-x_2^-$; il vincolo $\ge$ prende un surplus; l'uguaglianza ha termine noto negativo e si moltiplica per $-1$:
\[\min\ x_1-2x_2^++2x_2^-\ \ \st\ x_1+x_2^+-x_2^--s_1=2,\ \ -x_1+x_2^+-x_2^-=1,\ \ x_1,x_2^+,x_2^-,s_1\ge0.\]
\textbf{Canonica}: $\max\{-x_1+2x_2^+-2x_2^-:\ -x_1-x_2^++x_2^-\le-2,\ \ x_1-x_2^++x_2^-\le-1,\ \ -x_1+x_2^+-x_2^-\le1,\ \ x\ge0\}$ (il valore ottimo è l'opposto).
\end{esbox}

\section{Risoluzione grafica}

Ogni vincolo $a_i^Tx\le b_i$ è associato all'\textbf{iperpiano} $\{x\in\R^n:a_i^Tx=b_i\}$ (in $\R^2$ una retta) e definisce il \textbf{semispazio chiuso} ammissibile $\{x:a_i^Tx\le b_i\}$. La regione ammissibile è l'intersezione di tutti i semispazi.

Le \textbf{curve di livello} dell'obiettivo sono gli iperpiani $c^Tx=z$: rette parallele, ognuna associata a un valore $z$. Il \textbf{gradiente} $\nabla(c^Tx)=c$ è perpendicolare alle curve di livello e indica la direzione di crescita.

\textbf{Procedura} (2 variabili): disegna la regione; disegna una curva di livello; spostala parallelamente nella direzione di $c$ (per un $\min$: nella direzione $-c$) fino all'ultima posizione in cui tocca ancora la regione. Il valore di quella retta è l'ottimo, i punti di contatto sono le soluzioni ottime. Per il coltivatore: $3x_L+5x_P=0,30,42,44$; l'ultima che tocca è $z=44$ in $(8,4)$ (vedi la figura di \S\ref{sec:colt}). Il punto esatto si trova risolvendo il sistema dei due vincoli attivi: $x_L+x_P=12$, $x_L+2x_P=16$.

\textbf{Osservazione}: la soluzione ottima si trova sulla \emph{frontiera} della regione ammissibile.

\subsection{Quante soluzioni ottime può avere una PL?}

\begin{center}
\begin{tikzpicture}[scale=0.36,>=Stealth]
\begin{scope}
\fill[blue!12] (0,0)--(10,0)--(10,2)--(8,4)--(4,6)--(0,6)--cycle;
\draw[->] (-0.3,0)--(11.5,0); \draw[->] (0,-0.3)--(0,8);
\draw[thick] (0,0)--(10,0)--(10,2)--(8,4)--(4,6)--(0,6)--cycle;
\draw[dashed,red] (3,7)--(11.5,1.9); \fill[red] (8,4) circle (5pt);
\node[font=\scriptsize,align=center] at (5.5,-2) {(a) ottimo unico\\ $\max 3x_1+5x_2$};
\end{scope}
\begin{scope}[xshift=14cm]
\fill[blue!12] (0,0)--(10,0)--(10,2)--(8,4)--(4,6)--(0,6)--cycle;
\draw[->] (-0.3,0)--(11.5,0); \draw[->] (0,-0.3)--(0,8);
\draw[thick] (0,0)--(10,0)--(10,2)--(8,4)--(4,6)--(0,6)--cycle;
\draw[ultra thick,red] (8,4)--(10,2);
\node[font=\scriptsize,align=center] at (5.5,-2) {(b) infinite soluzioni ottime\\ $\max x_1+x_2$ (spigolo)};
\end{scope}
\begin{scope}[xshift=28cm]
\fill[blue!12] (0,0)--(11,0)--(11,5)--(0,5)--cycle;
\draw[->] (-0.3,0)--(11.5,0); \draw[->] (0,-0.3)--(0,8);
\draw[thick] (11,0)--(0,0)--(0,5)--(11,5);
\draw[->,thick,red] (4,2.5)--(8,2.5);
\node[font=\scriptsize,align=center] at (5.5,-2) {(c) illimitato\\ $\max x_1$, \ $0\le x_2\le5$, $x_1\ge0$};
\end{scope}
\end{tikzpicture}
\end{center}
\begin{itemize}[nosep]
\item \textbf{(a) Unica} soluzione ottima, in un vertice.
\item \textbf{(b) Infinite} soluzioni ottime: la curva di livello è parallela a un lato della regione; ottimi tutti i punti del segmento tra $(8,4)$ e $(10,2)$, valore 12. Gli estremi del segmento sono vertici.
\item \textbf{(c) Problema illimitato}: regione illimitata e obiettivo che cresce lungo una direzione in cui ci si può muovere all'infinito. Attenzione: \emph{regione illimitata non implica problema illimitato}: sulla stessa regione $\min x_1$ ha ottimo (segmento $x_1=0$) e $\max x_2$ ha ottimo $x_2=5$ (infiniti punti).
\item \textbf{Problema inammissibile}: $X=\emptyset$, nessuna soluzione (es.\ $x_1+x_2\le1$, $x_1+x_2\ge2$).
\end{itemize}
\begin{esame}
Un problema di PL è sempre in uno di questi casi: \textbf{inammissibile}, \textbf{illimitato}, oppure ha \textbf{ottimo finito}, raggiunto in 1 punto o in infiniti punti. \textbf{Mai esattamente 2} (o un numero finito $>1$) soluzioni ottime. Con variabili intere invece può succedere: $\max\{x_1+x_2: x_1+x_2\le1,\ x\in\{0,1\}^2\}$ ha esattamente 2 soluzioni ottime.
\end{esame}

\section{Geometria della PL}

\subsection{Combinazioni}
Dati $x^1,\dots,x^k\in\R^n$ e pesi $\lambda_1,\dots,\lambda_k$, il punto $y=\sum_i\lambda_ix^i$ è una combinazione
\begin{center}
\begin{tabular}{@{}lll@{}}
\toprule
\textbf{lineare} & $\lambda_i\in\R$ & in $\R^2$, di 2 vettori indipendenti: tutto il piano\\
\textbf{affine} & $\sum_i\lambda_i=1$ & di 2 punti: la retta che li contiene\\
\textbf{conica} & $\lambda_i\ge0$ & di 2 vettori: l'``angolo'' tra le due semirette dall'origine\\
\textbf{convessa} & $\lambda_i\ge0$, $\sum_i\lambda_i=1$ & di 2 punti: il segmento; di 3: il triangolo\\
\bottomrule
\end{tabular}
\end{center}

\begin{esbox}[esempi delle slide]
\begin{itemize}[nosep]
\item $(2,2)$ è combinazione convessa di $(4,0)$ e $(1,3)$: $\lambda(4,0)+(1-\lambda)(1,3)=(1+3\lambda,\ 3-3\lambda)=(2,2)$ per $\lambda=\tfrac13\in[0,1]$.
\item $(2,2)$ è combinazione convessa di $(1,1),(3,1),(2,3)$: dal sistema $\lambda_1+3\lambda_2+2\lambda_3=2$, $\lambda_1+\lambda_2+3\lambda_3=2$, $\lambda_1+\lambda_2+\lambda_3=1$ si ottiene $\lambda=(\tfrac14,\tfrac14,\tfrac12)\ge0$.
\item $(4,4)$ è combinazione conica di $(2,3)$ e $(3,1)$: $2a+3b=4$, $3a+b=4$ danno $a=\tfrac87$, $b=\tfrac47$, entrambi $\ge0$. (Non è convessa: $a+b=\tfrac{12}{7}\ne1$.)
\end{itemize}
Metodo: imposta il sistema lineare nei pesi e controlla i segni (e la somma).
\end{esbox}

\subsection{Insiemi convessi, coni, involucri}
\begin{defbox}[insieme convesso]
$S\subseteq\R^n$ è \textbf{convesso} se per ogni $x,y\in S$ e $\lambda\in[0,1]$ si ha $\lambda x+(1-\lambda)y\in S$ (contiene il segmento tra due suoi punti qualsiasi).
\end{defbox}
\begin{defbox}[cono]
$C\subseteq\R^n$ è un \textbf{cono} se $x\in C\Rightarrow \lambda x\in C$ per ogni $\lambda\ge0$. Un cono non è necessariamente convesso: l'unione dei due semiassi positivi $\{x\ge0: x_1x_2=0\}$ è un cono non convesso; il primo ortante $\R^2_+$ è un cono convesso.
\end{defbox}
\begin{defbox}[involucro convesso e conico]
L'\textbf{involucro (inviluppo) convesso} $\mathrm{conv}(K)$ è l'insieme di tutte le combinazioni convesse di elementi di $K$: il più piccolo convesso che contiene $K$. L'\textbf{involucro conico} $\mathrm{cono}(K)$ è l'insieme di tutte le combinazioni coniche: il più piccolo cono convesso che contiene $K$.
\end{defbox}
Esempio: $\mathrm{conv}\{(0,0),(2,0),(1,2)\}$ è il triangolo pieno con quei vertici, cioè $\{x: x_2\ge0,\ x_2\le2x_1,\ x_2\le4-2x_1\}$.

\subsection{Poliedri}
\textbf{Proprietà.}
\begin{itemize}[nosep]
\item L'intersezione di insiemi convessi (anche infiniti) è convessa.
\item Per $f$ lineare, $\{x: f(x)\le b\}$ è convesso per ogni $b$; $\{x\in\R^n: a^Tx\le b\}$ ($a\ne0$) è un \textbf{semispazio chiuso}.
\end{itemize}
\begin{defbox}[poliedro e politopo]
Un \textbf{poliedro} in $\R^n$ è l'intersezione di un numero \emph{finito} di semispazi chiusi; equivalentemente l'insieme delle soluzioni di un sistema di $m$ disequazioni lineari, $P=\{x\in\R^n: Ax\le b\}$. Un \textbf{politopo} è un poliedro limitato.
\end{defbox}
Quindi \textbf{la regione ammissibile di una PL è un poliedro} (convesso e chiuso); i vincoli $x\ge0$ sono semispazi come gli altri ($-x_j\le0$), le uguaglianze sono coppie di semispazi.

\emph{E con un vincolo non lineare?} In generale no: $\{x: x_1^2+x_2^2\le1\}$ è convesso ma non è un poliedro (servirebbero infiniti semispazi); $\{x: x_1^2+x_2^2\ge1\}$ non è nemmeno convesso.

\textbf{Coni e poliedri.} Non tutti i coni convessi sono poliedri: il cono di secondo ordine (``gelato'') $\{x\in\R^3: x_3\ge\sqrt{x_1^2+x_2^2}\}$ è un cono convesso ma non poliedrico.
\begin{prop}
$P$ è un \textbf{cono poliedrico} se e solo se esiste una matrice $Q$ tale che $P=\{x\in\R^n: Qx\le0\}$ (termini noti tutti nulli).
\end{prop}

\subsection{Vertici e direzioni di recessione}
\begin{defbox}[vertice]
$v\in P$ è un \textbf{vertice} (punto estremo) di $P$ se non si può scrivere come combinazione convessa \emph{stretta} di due punti distinti di $P$: se $v=\lambda x+(1-\lambda)y$ con $x,y\in P$ e $\lambda\in(0,1)$, allora $x=y=v$.
\end{defbox}
Caratterizzazione algebrica (utile per il simplesso): in $P=\{x\in\R^n:Ax\le b\}$, $v$ è un vertice se e solo se tra i vincoli \emph{attivi} in $v$ ($a_i^Tv=b_i$) ce ne sono $n$ linearmente indipendenti. In $\R^2$: un vertice è l'intersezione di due rette di vincolo non parallele. Il coltivatore ha 6 vertici: $(0,0),(10,0),(10,2),(8,4),(4,6),(0,6)$. Il punto $(12,0)$ è intersezione di due rette di vincolo ma non è ammissibile, quindi \emph{non} è un vertice.

\begin{defbox}[direzione di recessione]
$d\ne0$ è una \textbf{direzione di recessione} di $P$ se per ogni $x\in P$ e ogni $\lambda\ge0$ si ha $x+\lambda d\in P$ (da qualunque punto ci si può muovere all'infinito lungo $d$ restando in $P$). Per $P=\{x:Ax\le b\}\neq\emptyset$: $d$ è di recessione se e solo se $Ad\le0$. L'insieme $\{d: Ad\le0\}$ è il \emph{cono di recessione}.
\end{defbox}
Un politopo non ha direzioni di recessione. Esempio delle slide: $P=\{x\ge0,\ x_2\le5\}$: $Ad\le0$ dà $-d_1\le0$, $-d_2\le0$, $d_2\le0$, cioè $d=(d_1,0)$ con $d_1>0$: l'unica direzione (a meno di scala) è $(1,0)$.

\subsection{Teorema di decomposizione dei poliedri}
\begin{prop}[decomposizione, Minkowski--Weyl]
Ogni poliedro $P$ si può scrivere come
\[ P=\mathrm{conv}\{v^1,\dots,v^m\}+\mathrm{cono}\{e^1,\dots,e^p\}, \]
cioè ogni $x\in P$ è $x=\sum_{i}\lambda_iv^i+\sum_j\mu_je^j$ con $\lambda_i\ge0$, $\sum_i\lambda_i=1$, $\mu_j\ge0$: somma di un \textbf{politopo} e di un \textbf{cono} generato dalle direzioni di recessione. Se $P$ ha vertici, si possono prendere come $v^i$ esattamente i vertici. Vale anche il viceversa: un insieme di questa forma è un poliedro.
\end{prop}
(La ``somma'' è di Minkowski: $A+B=\{a+b: a\in A, b\in B\}$.) Il politopo o il cono possono essere banali: se $P$ è un politopo, $P=\mathrm{conv}(V)$ con $V$ insieme dei vertici.

\begin{esbox}[decomposizioni delle slide]
\begin{itemize}[nosep]
\item $P_1=\{x\in\R^2: x_1\ge1,\ x_2\ge1,\ x_1+x_2\ge3\}=\mathrm{conv}\{(1,2),(2,1)\}+\mathrm{cono}\{(1,0),(0,1)\}$.
\item $P_2=\{1\le x_1\le4,\ 1\le x_2\le3\}=\mathrm{conv}\{(1,1),(4,1),(4,3),(1,3)\}$ (politopo, cono $=\{0\}$).
\item $P_3=\{x_1\ge0,\ 0\le x_2\le1\}=\mathrm{conv}\{(0,0),(0,1)\}+\mathrm{cono}\{(1,0)\}$.
\item $P_4=\{0\le x_2\le1\}$ (striscia): \emph{nessun vertice} (contiene rette). $P_4=\mathrm{conv}\{(0,0),(0,1)\}+\mathrm{cono}\{(1,0),(-1,0)\}$: i punti $v^i$ esistono ma non sono vertici.
\end{itemize}
\end{esbox}
Un poliedro ha almeno un vertice se e solo se non contiene rette. Ogni poliedro della forma $\{x: Ax\le b, x\ge0\}$ non vuoto (quindi ogni PL in forma canonica o standard) ha vertici.

\section{Proprietà della soluzione ottima}

\begin{prop}[caratterizzazione elementare dell'ottimalità]
Dato $(P)\ \max\{c^Tx: x\in P\}$ con $P=\{x\in\R^n:Ax\le b\}$:
\begin{enumerate}[nosep]
\item se $c\ne0$, nessuna soluzione ottima è un punto interno di $P$;
\item se $(P)$ ha due soluzioni ottime, allora ne ha infinite;
\item ogni soluzione ottima locale è anche ottima globale.
\end{enumerate}
\end{prop}
\begin{dimo}
\emph{1.} Se $x^*$ è interno, esiste $\varepsilon>0$ tale che $x^*+\varepsilon c\in P$. Ma $c^T(x^*+\varepsilon c)=c^Tx^*+\varepsilon\|c\|^2>c^Tx^*$: $x^*$ non è ottima.
\emph{2.} Se $x',x''$ sono ottime con valore $z^*$, per ogni $\lambda\in[0,1]$ il punto $\lambda x'+(1-\lambda)x''$ sta in $P$ (convesso) e ha valore $\lambda z^*+(1-\lambda)z^*=z^*$ (linearità): tutto il segmento è ottimo.
\emph{3.} Sia $\bar x$ ottimo locale e supponiamo esista $y\in P$ con $c^Ty>c^T\bar x$. I punti $\bar x+\lambda(y-\bar x)$, $\lambda\in(0,1]$, stanno in $P$ e hanno valore $c^T\bar x+\lambda(c^Ty-c^T\bar x)>c^T\bar x$; per $\lambda$ piccolo sono vicini quanto si vuole a $\bar x$, contro la sua ottimalità locale. \hfill$\square$
\end{dimo}
Le proprietà 2 e 3 valgono per qualunque problema con regione convessa e obiettivo lineare (la 3 anche con obiettivo concavo da massimizzare); sono la ragione per cui la PL è ``facile'' rispetto alla PLI.

\begin{prop}[Teorema fondamentale della PL]
Sia $P=\mathrm{conv}\{v^1,\dots,v^m\}+\mathrm{cono}\{e^1,\dots,e^p\}$. Se il problema $(P)\ \max\{c^Tx: x\in P\}$ ha valore ottimo finito, allora esiste $k\in\{1,\dots,m\}$ tale che $v^k$ è una soluzione ottima.
\end{prop}
\begin{dimo}
\emph{Passo 1.} $c^Te^j\le0$ per ogni $j$. Altrimenti, preso un $\bar x\in P$, i punti $\bar x+te^j$ ($t\ge0$) stanno in $P$ e $c^T(\bar x+te^j)=c^T\bar x+t\,c^Te^j\to+\infty$: il problema sarebbe illimitato.
\emph{Passo 2.} Sia $k$ tale che $c^Tv^k=\max_ic^Tv^i$. Ogni $x\in P$ è $x=\sum_i\lambda_iv^i+\sum_j\mu_je^j$ con $\lambda\ge0$, $\sum\lambda_i=1$, $\mu\ge0$, quindi
\[ c^Tx=\sum_i\lambda_i\,c^Tv^i+\sum_j\mu_j\,c^Te^j\ \le\ \sum_i\lambda_i\,c^Tv^i\ \le\ \Big(\sum_i\lambda_i\Big)c^Tv^k=c^Tv^k. \]
Poiché $v^k\in P$, $v^k$ è ottimo. \hfill$\square$
\end{dimo}

\begin{oss}[conseguenze]
\begin{itemize}[nosep]
\item Se una PL (con vertici) ha ottimo finito, \textbf{esiste un vertice ottimo}. Basta cercare tra i vertici, che sono in numero finito: al più $\binom{m}{n}$ (scelte di $n$ vincoli attivi tra $m$).
\item Enumerare tutti i vertici è però esponenziale. Il \textbf{simplesso} si muove invece da un vertice a uno \emph{adiacente} che migliora l'obiettivo, e si ferma quando nessun vertice adiacente migliora (per la proprietà 3, un ottimo ``locale sui vertici'' è globale).
\item Il teorema dice che \emph{esiste} un vertice ottimo, non che tutti gli ottimi sono vertici (caso (b): anche i punti interni allo spigolo sono ottimi).
\item Con il Passo 1: una PL ammissibile è illimitata se e solo se esiste una direzione di recessione $e$ con $c^Te>0$.
\end{itemize}
\end{oss}

\begin{esame}
Da saper fare: portare un problema in forma standard e canonica; risolvere graficamente un problema in 2 variabili (anche di minimo) e riconoscere i casi ottimo unico / infiniti ottimi / illimitato / inammissibile; verificare se un punto è combinazione convessa o conica di altri; dire se un punto è un vertice; trovare le direzioni di recessione ($Ad\le0$); scrivere la decomposizione $\mathrm{conv}+\mathrm{cono}$ di un poliedro in $\R^2$; enunciare e dimostrare la caratterizzazione elementare e il teorema fondamentale.
\end{esame}

\chapter{Esercitazione 1: modelli (soluzioni)}\label{cap:es}

\emph{Soluzioni proposte per l'Esercitazione n.\,1 (Carello \& Pascoal). Molti esercizi ammettono formulazioni diverse ma equivalenti; qui si indica una formulazione corretta e, dove utile, il ragionamento e lo schema ricorrente. Per ogni esercizio: prova prima da solo, poi confronta. In aula (25/9) sono stati svolti gli esercizi \textbf{1, 6, 8 e 9}: sono quelli di riferimento per lo stile atteso.}

%------------------------------------------------
\section*{Es.\ 1 -- Miscelazione di benzine}\addcontentsline{toc}{section}{Es.\ 1 -- Miscelazione di benzine}
\textbf{Insiemi/parametri}: grezzi $I=\{1,2,3,4\}$ con disponibilità $a_i$ e costo $c_i$; benzine $J=\{A,B,C\}$ con prezzo $r_j$.
\textbf{Variabili}: $x_{ij}\ge0$ barili di grezzo $i$ usati nella benzina $j$. (La quantità prodotta di $j$ è $\sum_ix_{ij}$: non serve una variabile a parte.)
\begin{align*}
\max\ & \sum_{j\in J}r_j\sum_{i\in I}x_{ij}-\sum_{i\in I}c_i\sum_{j\in J}x_{ij}\\
\st\ & \sum_{j\in J}x_{ij}\le a_i && \forall i\in I\\
& x_{2A}\ge0.2\sum_{i}x_{iA},\qquad x_{3A}\le0.3\sum_ix_{iA} &&\text{(normale)}\\
& x_{3B}\ge0.4\sum_ix_{iB} &&\text{(super95)}\\
& x_{2C}\le0.5\sum_ix_{iC} &&\text{(super98)}\\
& x_{ij}\ge0.
\end{align*}
Schema: \emph{blending} (\S\ref{sec:blending}). Le percentuali si linearizzano moltiplicando per il totale della miscela.

%------------------------------------------------
\section*{Es.\ 2 -- Turni in ospedale}\addcontentsline{toc}{section}{Es.\ 2 -- Turni in ospedale}
Giorni $D=\{1,\dots,7\}$ (Lun$=1$), richiesta $r_k$.
\textbf{Idea chiave}: la variabile ``numero di infermieri in servizio il giorno $k$'' non basta, perché non garantisce i 5 giorni consecutivi. Si decide invece \textbf{quando inizia} il turno:
$x_j\in\Z_{\ge0}$ = numero di infermieri che iniziano i 5 giorni lavorativi il giorno $j$.
Il giorno $k$ lavorano quelli che hanno iniziato in $k,k-1,\dots,k-4$ (ciclicamente, modulo 7):
\[ \min\sum_{j\in D}x_j\qquad\st\ \sum_{j\in D_k}x_j\ge r_k\ \ \forall k\in D,\qquad x_j\in\Z_{\ge0}\ \ \forall j\in D, \]
dove $D_k=\{k,k-1,k-2,k-3,k-4\}$ con indici presi ciclicamente (modulo 7).
Ad esempio lunedì: $x_{\text{Gio}}+x_{\text{Ven}}+x_{\text{Sab}}+x_{\text{Dom}}+x_{\text{Lun}}\ge11$ (equivalentemente $\sum_jx_j-x_{\text{Mar}}-x_{\text{Mer}}\ge 11$).
\textbf{Soluzione}: ottimo intero 14 infermieri (es.\ $x=(6,3,0,4,1,0,0)$); il rilassamento continuo vale $13.\overline3$, quindi l'interezza conta davvero.

%------------------------------------------------
\section*{Es.\ 3 -- Posizionamento di ambulanze}\addcontentsline{toc}{section}{Es.\ 3 -- Posizionamento di ambulanze}
Punti di richiesta $I$, siti $J$, tempi $t_{ij}$; $a_{ij}=1$ se $t_{ij}\le8$ minuti (parametro calcolato dai dati). $y_j\in\{0,1\}$: ambulanza in $j$.
\[ \min\sum_{j\in J}y_j\qquad\st\ \sum_{j\in J}a_{ij}y_j\ge1\ \ \forall i\in I,\qquad y_j\in\{0,1\}. \]
Set covering. Nota: la soglia degli 8 minuti non va messa come vincolo sulle variabili, ma usata per \emph{pre-calcolare} $a_{ij}$ (oppure si somma solo su $j$ con $t_{ij}\le 8$).

%------------------------------------------------
\section*{Es.\ 4 -- Gestione dei rifiuti}\addcontentsline{toc}{section}{Es.\ 4 -- Gestione dei rifiuti}
\textbf{Variabili} (tonnellate, $\ge0$): $x_{ij}$ da sito $i$ a discarica $j$; $u_{ih}$ da sito $i$ a centro $h$; $v_{hj}$ scarti da centro $h$ a discarica $j$; $z_h$ materiale da $h$ al riciclaggio $r$; $y_h\in\{0,1\}$ costruisco il centro $h$.
\begin{align*}
\min\ & \sum_hf_hy_h+\sum_h p_h\sum_iu_{ih}+\sum_jb_j\Big(\sum_ix_{ij}+\sum_hv_{hj}\Big)+{}\\
&\quad\sum_{i,j}c_{ij}x_{ij}+\sum_{i,h}c_{ih}u_{ih}+\sum_{h,j}c_{hj}v_{hj}+\sum_hc_{hr}z_h-g\sum_hz_h\\
\st\ & \sum_jx_{ij}+\sum_hu_{ih}=t_i && \forall i\in I\\
& \sum_iu_{ih}\le q_hy_h && \forall h\in H\\
& z_h=0.8\sum_iu_{ih},\qquad \sum_jv_{hj}=0.2\sum_iu_{ih} && \forall h\in H\\
& \sum_ix_{ij}+\sum_hv_{hj}\le d_j && \forall j\in D
\end{align*}
Vincoli nell'ordine: tutti i rifiuti vanno smaltiti; capacità del centro (solo se costruito); bilanci del centro (80\% riciclato, 20\% in discarica), capacità delle discariche. Qui $p_h$ è inteso come costo di lavorazione per tonnellata trattata. Schema: \emph{flussi con bilanciamento} + \emph{costo fisso} con capacità ($\sum u\le q_hy_h$ fa da big-$M$ con $M=q_h$).

%------------------------------------------------
\section*{Es.\ 5 -- Linea ferroviaria}\addcontentsline{toc}{section}{Es.\ 5 -- Linea ferroviaria}
Il segmento $[0,L]$ contiene infiniti punti: bisogna \textbf{discretizzare}. L'antenna $i$ copre $[d_i-5,d_i+5]$. Siano $0=\pi_0<\pi_1<\dots<\pi_K=L$ i punti distinti dell'insieme $\{0,L\}\cup\{d_i\pm5\}\cap[0,L]$ ordinati. Dentro ogni intervallo elementare $(\pi_k,\pi_{k+1})$ nessuna antenna ``inizia'' o ``finisce'', quindi l'insieme delle antenne che lo coprono è lo stesso per tutti i suoi punti: basta controllare il punto medio $m_k$. Parametro $a_{ik}=1$ se $|d_i-m_k|\le5$.
\[ \min\sum_{i=1}^nx_i\qquad\st\ \sum_{i=1}^na_{ik}x_i\ge1\ \ \forall k=0,\dots,K-1,\qquad x_i\in\{0,1\}. \]
(Una classe di variabili, una classe di vincoli, come suggerito. Gli estremi $\pi_k$ sono coperti perché gli intervalli di copertura sono chiusi.)
\textbf{Variante} con capacità: $\sum_ia_{ik}c_ix_i\ge C\ \ \forall k$.

%------------------------------------------------
\section*{Es.\ 6 -- Manutenzione stradale}\addcontentsline{toc}{section}{Es.\ 6 -- Manutenzione stradale}
Grafo $G=(V,E)$ delle città e strade, $x_i\in\{0,1\}$ = ufficio in $i$. $E_L=\{\{i,j\}\in E: l_{ij}>L\}$.
\[ \min\sum_{i\in V}x_i\quad\st\ x_i+x_j\ge1\ \ \forall\{i,j\}\in E\setminus E_L;\qquad x_i+x_j\ge 2\ \ (\text{cioè }x_i=x_j=1)\ \ \forall\{i,j\}\in E_L. \]
È un \emph{vertex cover} con alcuni nodi forzati.
\textbf{Variante} (ogni ufficio gestisce al più 3 strade): serve sapere \emph{chi} gestisce ogni strada. $w_{ei}\in\{0,1\}$ per $e\in E$, $i\in e$: la strada $e$ è gestita dall'ufficio in $i$.
\[ \sum_{i\in e}w_{ei}\ge1\ \ \forall e\in E;\qquad w_{ei}\le x_i;\qquad \sum_{e\ni i}w_{ei}\le 3x_i\ \ \forall i\in V;\qquad w_{ei}=1\ \ \forall e\in E_L,\ i\in e. \]
(L'ultimo vincolo interpreta ``incluse quelle lunghe'': una strada lunga è gestita da entrambi gli uffici e conta per entrambi; esso implica anche $x_i=1$ per gli estremi delle strade lunghe.)

%------------------------------------------------
\section*{Es.\ 7 -- Carte da gioco}\addcontentsline{toc}{section}{Es.\ 7 -- Carte da gioco}
Carte $C=\{1,\dots,n\}$, simboli candidati $S=\{1,\dots,\bar S\}$ con $\bar S$ limite superiore (es.\ $\bar S=nk$).
Variabili: $x_{cs}\in\{0,1\}$ simbolo $s$ sulla carta $c$; $u_s\in\{0,1\}$ simbolo $s$ usato; $z_{cc's}\in\{0,1\}$ ($c<c'$) simbolo $s$ comune a $c$ e $c'$ (AND linearizzato).
\begin{align*}
\min\ & q=\sum_{s\in S}u_s\\
\st\ & \sum_sx_{cs}=k && \forall c\\
& x_{cs}\le u_s && \forall c,s\\
& z_{cc's}\le x_{cs},\quad z_{cc's}\le x_{c's},\quad z_{cc's}\ge x_{cs}+x_{c's}-1 && \forall c<c',\ s\\
& \sum_sz_{cc's}=1 && \forall c<c'.
\end{align*}
\textbf{Variante} ($\bar q$ simboli dati, massimizzare le carte): carte potenziali $c\in\{1,\dots,\bar N\}$ con $v_c\in\{0,1\}$ = carta creata; $\sum_{s=1}^{\bar q}x_{cs}=kv_c$; $\sum_sz_{cc's}\ge v_c+v_{c'}-1$ e $\sum_sz_{cc's}\le 1$; $\max\sum_cv_c$. (Se una carta non è creata non ha simboli, quindi le $z$ che la coinvolgono valgono 0.)

%------------------------------------------------
\section*{Es.\ 8 -- Il pastificio}\addcontentsline{toc}{section}{Es.\ 8 -- Il pastificio}
$T=\{1,\dots,|T|\}$; $x_t\ge0$ quintali prodotti; $s_t\ge0$ scorta a fine giorno ($s_0=0$); $y_t\in\{0,1\}$ macchine attive.
\begin{align*}
\min\ & \sum_{t\in T}(Sy_t+c_tx_t+gs_t)\\
\st\ & s_{t-1}+x_t=d_t+s_t && \forall t\in T\\
& qy_t\le x_t\le Qy_t && \forall t\in T\\
& 0\le s_t\le B && \forall t\in T,\qquad s_0=0.
\end{align*}
Lot sizing: bilanciamento + lotto minimo + costo fisso.

%------------------------------------------------
\section*{Es.\ 9 -- Dirigente scolastico}\addcontentsline{toc}{section}{Es.\ 9 -- Dirigente scolastico}
$y_j\in\{0,1\}$ attivo il corso $j$.
\[ \min\sum_{j\in C}g_jy_j\qquad\st\ \sum_{j\in f_i\cup s_i}y_j\ge1\ \ \forall i\in I. \]
\textbf{Variante}: $z_i\in\{0,1\}$ = lo studente $i$ ha almeno una prima scelta attiva. $z_i\le\sum_{j\in f_i}y_j\ \forall i$ e $\sum_iz_i\ge0.5|I|$. Basta la direzione ``$z_i=1\Rightarrow$ almeno una prima scelta'': il vincolo sul 50\% spinge le $z$ verso 1, non serve impedire che valgano 0.

%------------------------------------------------
\section*{Es.\ 10 -- Doni natalizi}\addcontentsline{toc}{section}{Es.\ 10 -- Doni natalizi}
Pacchi potenziali $K=\{1,\dots,\bar K\}$ con $\bar K=\lfloor\sum_iv_i/V\rfloor$ (o semplicemente $n$). $x_{ik}\in\{0,1\}$ oggetto $i$ nel pacco $k$; $y_k\in\{0,1\}$ pacco $k$ confezionato.
\[ \max\sum_ky_k\quad\st\ \sum_kx_{ik}\le1\ \forall i;\qquad \sum_iv_ix_{ik}\ge Vy_k\ \forall k. \]
(Un pacco non ``attivo'' può contenere oggetti senza problemi: semplicemente non conta.)
\textbf{Variante}: $x_{ik}+x_{jk}\le1\ \ \forall(i,j)\in A,\ \forall k\in K$.

%------------------------------------------------
\section*{Es.\ 11 -- Setup di robot}\addcontentsline{toc}{section}{Es.\ 11 -- Setup di robot}
$x_{ij}\ge0$ unità del prodotto $i$ sul robot $j$ (intere se i prodotti sono indivisibili); $y_{ij}\in\{0,1\}$ setup di $j$ per $i$; $z$ produzione della tipologia meno prodotta.
\begin{align*}
\max\ & z\\
\st\ & z\le\sum_jx_{ij} && \forall i &&\text{(max-min)}\\
& \sum_ia_{ij}x_{ij}\le b_j && \forall j\\
& x_{ij}\le \tfrac{b_j}{a_{ij}}\,y_{ij} && \forall i,j &&\text{(big-$M$ con $M=b_j/a_{ij}$)}\\
& \sum_ir_i\sum_jx_{ij}-\sum_{i,j}c_{ij}y_{ij}\ge B &&&&\text{(ritorno minimo)}
\end{align*}

%------------------------------------------------
\section*{Es.\ 12 -- Azienda dolciaria}\addcontentsline{toc}{section}{Es.\ 12 -- Azienda dolciaria}
Mesi $I=\{1,\dots,|I|\}$. Variabili: $p_{ik}\in\Z_{\ge0}$ cioccolatini $k$ prodotti nel mese $i$; $s_{ik}\ge0$ scorta a fine mese ($s_{0k}=0$); $o_i\ge0$ minuti di straordinario; $w_i\in\{0,1\}$ straordinari attivati; $n_j\in\Z_{\ge0}$ confezioni di materia prima $j$ acquistate.
\begin{align*}
\min\ & \sum_{i,k}c_kp_{ik}+\sum_jp_jn_j+\sum_if_iw_i+\sum_{i,k}g_ks_{ik}\\
\st\ & s_{i-1,k}+p_{ik}=d_{ik}+s_{ik} && \forall i,k\\
& \sum_kh_kp_{ik}\le Q+o_i && \forall i\\
& o_i\le 60S\,w_i && \forall i &&(S\text{ in ore}, Q\text{ in minuti})\\
& v_jn_j\ge\sum_{i,k}a_{kj}p_{ik} && \forall j &&\text{(acquisto a inizio periodo)}\\
& \sum_iw_i\le b,\qquad w_i+w_{i+1}\le1 && \forall i<|I|
\end{align*}

%------------------------------------------------
\section*{Es.\ 13 -- Localizzazione di discariche}\addcontentsline{toc}{section}{Es.\ 13 -- Localizzazione di discariche}
Aree $A=\{1,\dots,10\}$, città $K=\{1,\dots,7\}$, $N_c\subseteq A$ aree vicine alla città $c$ (dalla tabella). $y_a\in\{0,1\}$ costruisco in $a$.
\begin{align*}
\max\ & \sum_aw_ay_a\\
\st\ & \sum_{a\in N_c}y_a\le1 && \forall c\in K\\
& y_2\le y_7;\qquad y_1+y_2\le1;\qquad y_1+y_4\le1\\
& \sum_ag_ay_a\le25.
\end{align*}
Con i dati della tabella l'ottimo è $\{4,9\}$ con capacità 1290 (costo 12). Il vincolo di vicinanza è molto restrittivo: le città 1, 2, 4 e 7 ``vedono'' molte aree.

%------------------------------------------------
\section*{Es.\ 14 -- Matrimonio}\addcontentsline{toc}{section}{Es.\ 14 -- Matrimonio}
Invitati $P$, gruppi $G_1,\dots,G_k$, tavoli potenziali $T$ ($|T|=\lceil|P|/p\rceil$ o più). Variabili: $x_{it}\in\{0,1\}$ invitato $i$ al tavolo $t$; $y_t\in\{0,1\}$ tavolo usato; $u_{gt}\in\{0,1\}$ gruppo $g$ presente al tavolo $t$.
\begin{align*}
\min\ & \sum_ty_t\\
\st\ & \sum_tx_{it}=1 && \forall i\\
& \sum_ix_{it}\le p\,y_t && \forall t\\
& u_{gt}\le\sum_{i\in G_g}x_{it} && \forall g,t\\
& \sum_gu_{gt}\ge q\,y_t && \forall t &&\text{(eterogeneità)}\\
& \sum_{j\in G_g\setminus\{i\}}x_{jt}\ge r\,x_{it} && \forall g,\ i\in G_g,\ t &&\text{(almeno $r$ compagni)}
\end{align*}
\textbf{Variante}: $\min z$ con $z\ge\sum_{i\in G_g}x_{it}\ \forall g,t$ (min-max), eliminando l'obiettivo sul numero di tavoli (i vincoli su $y$ restano per definire i tavoli usati).

%------------------------------------------------
\section*{Es.\ 15 -- Instradamento di pacchetti}\addcontentsline{toc}{section}{Es.\ 15 -- Instradamento di pacchetti}
$x_{i1},x_{i2}\in\{0,1\}$: pacchetto $i$ sul link 1 / 2; $r_i$ risorse, $c_i$ costo sul link 1.
\[ \min\sum_i\big(c_ix_{i1}+1.3c_ix_{i2}\big)\quad\st\ x_{i1}+x_{i2}=1\ \forall i;\quad \sum_ir_ix_{i1}\le1;\quad\sum_ir_ix_{i2}\le2. \]
\textbf{Soluzione}: sul link 1 vanno i pacchetti $\{4,6,8,10\}$ (0.6+0.2+0.1+0.1 = 1 Mbps, costo 1400), il resto sul link 2 (1.7 Mbps): costo totale 3610. Ragionamento: costo $=1.3\sum_ic_i-0.3\sum_{i\in\text{link 1}}c_i$, quindi si risolve uno \emph{zaino} che massimizza il costo messo sul link 1.
\textbf{Variante} con $m$ link: $x_{il}$, $\sum_lx_{il}=1\ \forall i$, $\sum_ir_ix_{il}\le k_l\ \forall l$, costo $c_{il}$.

%------------------------------------------------
\section*{Es.\ 16 -- Fabbrica di bulloni}\addcontentsline{toc}{section}{Es.\ 16 -- Fabbrica di bulloni}
$x_i\in[p,q]$ produzione del mese $i$; $s_i\ge0$ giacenza a fine mese ($s_0=m_0$); $y_i\in\{0,1\}$ magazzino raddoppiato nel mese $i$.
\[ \min\sum_i(r\,s_i+f\,y_i)\quad\st\ s_{i-1}+x_i=d_i+s_i,\quad s_i\le M+My_i,\quad p\le x_i\le q,\quad s_i\ge0. \]
(Se il costo $r$ si paga sulla giacenza media o iniziale, cambia solo l'indice nell'obiettivo.)

Sia $D_t=\sum_{i\le t}d_i$. La giacenza a fine mese $t$ è $s_t=m_0+\sum_{i\le t}x_i-D_t\in[m_0+tp-D_t,\ m_0+tq-D_t]$.
\begin{itemize}[nosep]
\item \textbf{Condizione sufficiente di inammissibilità}: esiste $t$ con $m_0+tq<D_t$ (anche producendo al massimo non si soddisfa la domanda) oppure $m_0+tp-D_t>2M$ (anche producendo al minimo il magazzino esplode). Caso più semplice: $d_1>m_0+q$.
\item \textbf{Condizione necessaria di ammissibilità}: la negazione della precedente, $m_0+tq\ge D_t$ e $m_0+tp-D_t\le2M$ per ogni $t$ (in particolare $m_0+nq\ge D_n$).
\end{itemize}

%------------------------------------------------
\section*{Es.\ 17 -- Rete di distribuzione}\addcontentsline{toc}{section}{Es.\ 17 -- Rete di distribuzione}
$y_j\in\{0,1\}$ costruisco il centro $j$; $x_{ij}\in\{0,1\}$ negozio $i$ collegato a $j$.
\[ \min\sum_jc_jy_j+\sum_{i,j}\alpha d_{ij}x_{ij}\quad\st\ \sum_jx_{ij}\ge1\ \forall i;\quad 2y_j\le\sum_ix_{ij}\le k\,y_j\ \forall j. \]
Il vincolo destro impone anche $x_{ij}=0$ se $y_j=0$.
\textbf{Variante}: nuova variabile $z$ e nuovi vincoli $z\ge d_{ij}x_{ij}\ \forall i,j$; obiettivo $\min z$.

%------------------------------------------------
\section*{Es.\ 18 -- Ambulanze}\addcontentsline{toc}{section}{Es.\ 18 -- Ambulanze}
$x_{ij}\in\{0,1\}$ città $i$ assegnata all'ospedale $j$; livello di servizio di $j$: $\sum_it_{ij}x_{ij}$.
\begin{enumerate}[nosep]
\item $\min z$ con $z\ge\sum_it_{ij}x_{ij}\ \forall j$ e $\sum_jx_{ij}=1\ \forall i$.
\item $\min\sum_j\sum_it_{ij}x_{ij}$ con gli stessi vincoli: il problema si separa per città e l'ottimo assegna ogni città all'\emph{ospedale più vicino}.
\item Nella figura (coordinate sulla griglia): città $C_1=(3,4)$, $C_2=(1,2)$, $C_3=(1,1)$; ospedali $H_1=(4,4)$, $H_2=(2,2)$, $H_3=(3,1)$.
\emph{Min-max}: $C_1\to H_1$, $C_2\to H_2$, $C_3\to H_3$, livelli $1,1,2$: \textbf{obiettivo 2}.
\emph{Totale}: ciascuna al più vicino, $C_1\to H_1$, $C_2\to H_2$, $C_3\to H_2$ (distanza $\sqrt2$): \textbf{obiettivo $2+\sqrt2\approx3.41$} (e un livello massimo di $1+\sqrt2\approx2.41$ per $H_2$).
\end{enumerate}
Morale: il min-max ``bilancia'', il min-somma ``efficienta''.

%------------------------------------------------
\section*{Es.\ 19 -- Medici di base}\addcontentsline{toc}{section}{Es.\ 19 -- Medici di base}
$x_{ij}\in\{0,1\}$ per $j\in m_i$.
\begin{enumerate}[nosep]
\item $\max\sum_i\sum_{j\in m_i}r_{ij}x_{ij}$ con $\sum_{j\in m_i}x_{ij}=1\ \forall i$ e $\sum_{i:j\in m_i}w_ix_{ij}\le W\ \forall j$.
\item $\sum_j\sum_i\alpha r_{ij}x_{ij}=\alpha\cdot$ (obiettivo del punto 1): stessi vincoli, obiettivo moltiplicato per una costante positiva $\Rightarrow$ \textbf{stesse soluzioni ottime}. La soluzione non è necessariamente diversa (sommare per medici o per pazienti è la stessa doppia somma).
\item $\max z$ con $z\le\sum_{j\in m_i}r_{ij}x_{ij}\ \forall i$.
\item $\max z$ con $z\le\sum_i\alpha r_{ij}x_{ij}\ \forall j$: ora il minimo è preso sui \emph{medici} (somme per colonna) e non sui pazienti. Il $\min$ non commuta con la somma, quindi in generale la soluzione ottima è \textbf{diversa}.
\end{enumerate}

%------------------------------------------------
\section*{Es.\ 20 -- Rete stradale}\addcontentsline{toc}{section}{Es.\ 20 -- Rete stradale}
$x_i\in\{0,1\}$ rotonda nell'incrocio $i$; $E_5=\{\{i,j\}: t_{ij}\ge5\}$.
\begin{enumerate}[nosep]
\item $\min\sum_ix_i$ con $x_i+x_j\ge1\ \forall\{i,j\}\in E_5$ (vertex cover sui lati ``trafficati'').
Lati con $t\ge5$: $ha(10),ab(6),ae(5),ad(7),dg(12),hb(8),bc(5),ef(5),fg(6)$. Ottimo: \textbf{4 rotonde}, es.\ $\{a,b,e,g\}$ oppure $\{a,b,d,f\}$.
\item Variante: $x_i=x_j=1$ (cioè $x_i+x_j\ge2$) per ogni lato con $t_{ij}\ge8$: $ha(10),hb(8),dg(12)$. Ottimo $\{a,b,d,g,h\}$, 5 rotonde.
\end{enumerate}

%------------------------------------------------
\section*{Es.\ 21 -- Baguette}\addcontentsline{toc}{section}{Es.\ 21 -- Baguette}
$n_{1i},n_{2i}\in\Z_{\ge0}$: lotti da 20 e da 50 il giorno $i$; $z_{1i},z_{2i}\in\{0,1\}$: si produce quel tipo di lotto.
\begin{align*}
\min\ & \sum_i(c_1z_{1i}+c_2z_{2i}+d_1n_{1i}+d_2n_{2i})\\
\st\ & 20n_{1i}+50n_{2i}\ge r_i && \forall i\\
& n_{1i}\le\lceil r_i/20\rceil z_{1i},\qquad n_{2i}\le\lceil r_i/50\rceil z_{2i} && \forall i
\end{align*}
(Nessun magazzino: la baguette si vende in giornata. Le $M$ sono i numeri massimi sensati di lotti.)
\begin{enumerate}[nosep,start=2]
\item $z_{1i}+z_{2i}\le1\ \ \forall i$.
\item $z_{2i}+z_{2,i+1}+z_{2,i+2}+z_{2,i+3}\le3\ \ \forall i=1,\dots,7$ (su 4 giorni consecutivi al più 3 di produzione da 50).
\end{enumerate}

%------------------------------------------------
\section*{Es.\ 22 -- Asili nido}\addcontentsline{toc}{section}{Es.\ 22 -- Asili nido}
$y_j\in\{0,1\}$ nido nel paese $j$; $x_{ij}\in\{0,1\}$ bambini di $i$ assegnati a $j$, definite solo per $t_{ij}\le\Delta$ (insieme $A$); $e_j\in\Z_{\ge0}$ educatori in $j$.
\begin{align*}
\min\ & \sum_{(i,j)\in A,\ i\ne j}c_{ij}x_{ij}+g\sum_je_j\\
\st\ & \sum_{j:(i,j)\in A}x_{ij}=1 && \forall i\\
& x_{ij}\le y_j && \forall(i,j)\in A\\
& \sum_jy_j\le p\\
& q\,e_j\ge\sum_{i:(i,j)\in A}b_ix_{ij} && \forall j &&(e_j\ge \text{bambini}/q\text{, intero})
\end{align*}
\textbf{Variante}: $\sum_ib_ix_{ij}\ge r\,y_j\ \forall j$; $y_j+y_l\le1$ per ogni coppia $j\neq l$ con $t_{jl}<\delta$.

%------------------------------------------------
\section*{Es.\ 23 -- Problema su grafo}\addcontentsline{toc}{section}{Es.\ 23 -- Problema su grafo}
\begin{enumerate}[nosep]
\item \textbf{Insieme indipendente massimo}: $\max\sum_ix_i$, $x_i+x_j\le1\ \forall\{i,j\}\in E$.
\item \textbf{Vertex cover minimo}: $\min\sum_ix_i$, $x_i+x_j\ge1\ \forall\{i,j\}\in E$.
\end{enumerate}
Relazione: $S$ è indipendente se e solo se $V\setminus S$ è un vertex cover (sostituisci $x_i\to1-x_i$). Quindi $|S^*|+|T^*|=|V|$.
Figura 1 (casa: pentagono con tetto e diagonale): indipendente massimo di cardinalità 2 (es.\ il vertice in alto e quello in basso a sinistra); vertex cover minimo di cardinalità 3 (i tre restanti). Figura 2 (pentagono con tutte le diagonali $=K_5$): indipendente massimo 1 (un vertice qualsiasi), vertex cover minimo 4.

\chapter[Allenamento quesiti brevi]{Allenamento per i quesiti brevi (Vero/Falso)}\label{cap:quiz}

Stesso formato dei quesiti brevi: per ogni affermazione rispondi V o F, in circa 30 secondi ciascuna. Le soluzioni motivate sono in fondo al capitolo. Le domande Q1--Q6 riguardano i Capitoli 1--3 (primo quesito, 6/10); Q7--Q8 la PL (Capitolo~\ref{cap:pl}).

\newcommand{\aff}[1]{\item[#1] }
\begin{description}[leftmargin=1.2cm,style=nextline,itemsep=1pt]
\item[\textbf{Q1}] \textbf{Definizioni.}
\begin{description}[leftmargin=1.2cm,itemsep=0pt]
\aff{1.1} Un problema di ottimizzazione può avere due soluzioni ottime con valori diversi della funzione obiettivo.
\aff{1.2} Se la regione ammissibile è vuota, il problema si dice illimitato.
\aff{1.3} $\min\{f(x):x\in X\}=-\max\{-f(x):x\in X\}$.
\aff{1.4} Il problema $\min\{x: x>0\}$ ha soluzione ottima $x=0$.
\end{description}

\item[\textbf{Q2}] \textbf{Classi di modelli.}
\begin{description}[leftmargin=1.2cm,itemsep=0pt]
\aff{2.1} Un modello con vincoli e obiettivo lineari, alcune variabili continue e alcune binarie è un modello di PLMI.
\aff{2.2} Il vincolo $x_1x_2\le 4$, con $x_1,x_2$ continue, è lineare.
\aff{2.3} Il vincolo $\frac{x_1}{x_1+x_2}\ge 0.3$ (con $x_1+x_2>0$) si può riscrivere in modo equivalente come vincolo lineare.
\aff{2.4} Il valore ottimo del rilassamento continuo di un problema di massimo in PLI è un limite superiore del suo valore ottimo.
\end{description}

\item[\textbf{Q3}] \textbf{Vincoli logici} ($y_a,y_b,y_c\in\{0,1\}$).
\begin{description}[leftmargin=1.2cm,itemsep=0pt]
\aff{3.1} ``Se si sceglie $a$ allora si sceglie $b$'' si scrive $y_b\le y_a$.
\aff{3.2} ``Si sceglie almeno una tra $a$ e $b$'' si scrive $y_a+y_b=1$.
\aff{3.3} ``Se si scelgono sia $a$ sia $b$, si sceglie $c$'' si scrive $y_a+y_b-1\le y_c$.
\aff{3.4} ``Se si sceglie almeno una tra $a$ e $b$, si sceglie $c$'' si può scrivere $y_a+y_b\le 2y_c$.
\end{description}

\item[\textbf{Q4}] \textbf{Big-$M$ e lotto minimo} ($x\ge0$ continua, $y\in\{0,1\}$).
\begin{description}[leftmargin=1.2cm,itemsep=0pt]
\aff{4.1} Il vincolo $x\le My$ impone che, se $y=1$, allora $x>0$.
\aff{4.2} Con $qy\le x\le My$ ($0<q<M$), $x$ può valere 0 oppure qualunque valore in $[q,M]$.
\aff{4.3} Scegliere $M$ molto più grande del necessario rende il modello errato (esclude soluzioni ammissibili).
\aff{4.4} Per attivare il vincolo $g(x)\le b$ solo quando $y=1$ si può scrivere $g(x)\le b+M(1-y)$, con $M\ge\max g(x)-b$.
\end{description}

\item[\textbf{Q5}] \textbf{Obiettivi e variabili.}
\begin{description}[leftmargin=1.2cm,itemsep=0pt]
\aff{5.1} $\min\max_i f_i(x)$, con $f_i$ lineari, equivale a $\min z$ con $z\ge f_i(x)$ per ogni $i$.
\aff{5.2} $\max\max_i f_i(x)$ equivale a $\max z$ con $z\ge f_i(x)$ per ogni $i$.
\aff{5.3} Una variabile libera $x$ si può sostituire con $x^+-x^-$, $x^+,x^-\ge0$.
\aff{5.4} Per far assumere a una variabile solo i valori $\{d_1,\dots,d_r\}$ servono $r$ variabili binarie con somma pari a 1.
\end{description}

\item[\textbf{Q6}] \textbf{Complessità.}
\begin{description}[leftmargin=1.2cm,itemsep=0pt]
\aff{6.1} $\NP$ significa ``non polinomiale''.
\aff{6.2} $\Pc\subseteq\NP$.
\aff{6.3} Se $P_1\red P_2$ e $P_2\in\Pc$, allora $P_1\in\Pc$.
\aff{6.4} Per dimostrare che un problema $P$ è $\NP$-completo basta mostrare che $P\red$ SAT.
\aff{6.5} Un algoritmo di complessità $O(NB)$ per lo zaino ($B$ capacità) è polinomiale nella dimensione dell'istanza.
\aff{6.6} La dimensione di un'istanza è il numero di bit necessari a codificarla.
\aff{6.7} Lo zaino (versione di riconoscimento) è in $\NP$ perché dato un sottoinsieme di oggetti si verifica in tempo polinomiale che pesi $\le B$ e abbia profitto $\ge k$.
\aff{6.8} La complessità di un algoritmo si misura nel caso medio.
\end{description}

\item[\textbf{Q7}] \textbf{Forme e soluzione grafica della PL.}
\begin{description}[leftmargin=1.2cm,itemsep=0pt]
\aff{7.1} Un problema di PL può avere esattamente due soluzioni ottime.
\aff{7.2} Se la regione ammissibile di una PL è illimitata, il problema è illimitato.
\aff{7.3} Il vincolo $a^Tx\ge b$ si porta in forma standard con una variabile di surplus: $a^Tx-s=b$, $s\ge0$.
\aff{7.4} In un problema di minimo, nella risoluzione grafica ci si muove nella direzione del gradiente dell'obiettivo.
\end{description}

\item[\textbf{Q8}] \textbf{Geometria della PL.}
\begin{description}[leftmargin=1.2cm,itemsep=0pt]
\aff{8.1} Il punto $(2,2)$ è combinazione convessa di $(4,0)$ e $(1,3)$.
\aff{8.2} Ogni cono è convesso.
\aff{8.3} L'insieme $\{x\in\R^2: x_1^2+x_2^2\le1\}$ è un poliedro.
\aff{8.4} Se una PL ha ottimo finito e la regione ammissibile ha vertici, almeno un vertice è ottimo.
\aff{8.5} Una soluzione ottima locale di una PL è anche globale.
\aff{8.6} La striscia $\{x\in\R^2: 0\le x_2\le1\}$ ha due vertici.
\end{description}
\end{description}

\section*{Soluzioni}
\begin{description}[leftmargin=1.2cm,itemsep=1pt,font=\normalfont\bfseries]
\item[1.1 F] Le soluzioni ottime possono essere più d'una, ma hanno tutte lo stesso valore (il valore ottimo è unico).
\item[1.2 F] Regione vuota $\Rightarrow$ problema \emph{inammissibile}. Illimitato: l'obiettivo migliora senza limite.
\item[1.3 V] E le soluzioni ottime coincidono.
\item[1.4 F] $0$ non è ammissibile e ogni $x>0$ è battuto da $x/2$: nessun ottimo. Per questo non si usano disuguaglianze strette.
\item[2.1 V] Mista intera = alcune variabili intere (le binarie lo sono), altre continue.
\item[2.2 F] Prodotto di variabili: non lineare.
\item[2.3 V] $x_1\ge0.3(x_1+x_2)$, moltiplicando per il denominatore positivo (vincolo di blending).
\item[2.4 V] Il rilassamento ha più soluzioni ammissibili; per un massimo dà un limite superiore.
\item[3.1 F] È $y_a\le y_b$: se $y_a=1$ deve essere $y_b=1$. $y_b\le y_a$ esprime ``se $b$ allora $a$''.
\item[3.2 F] Quello è ``esattamente una''. ``Almeno una'' è $y_a+y_b\ge1$.
\item[3.3 V] Se $y_a=y_b=1$ il lato sinistro vale 1, quindi $y_c=1$; negli altri casi vale $\le0$.
\item[3.4 V] Se almeno una vale 1, il lato sinistro è $\ge1$ e serve $y_c=1$; il 2 permette il caso $y_a=y_b=1$.
\item[4.1 F] Impone solo $x>0\Rightarrow y=1$ (cioè $y=0\Rightarrow x=0$). Con $y=1$, $x=0$ è ancora ammesso.
\item[4.2 V] Insieme $\{0\}\cup[q,M]$, non convesso: richiede la binaria.
\item[4.3 F] Il modello resta corretto se $M$ è un limite valido; un $M$ eccessivo indebolisce il rilassamento continuo e crea problemi numerici, ma non esclude soluzioni.
\item[4.4 V] Con $y=0$ il vincolo diventa ridondante se $b+M$ supera il massimo di $g$.
\item[5.1 V] All'ottimo $z$ è uguale al massimo (altrimenti si potrebbe diminuire).
\item[5.2 F] Con $z\ge f_i$ e $\max z$, $z$ cresce all'infinito. Max-max (e min-min) richiedono variabili binarie.
\item[5.3 V] È la trasformazione usata per la forma standard.
\item[5.4 V] $\sum_k y_k=1$ e $x=\sum_kd_ky_k$.
\item[6.1 F] Non-deterministico Polinomiale.
\item[6.2 V] Chi sa risolvere in tempo polinomiale sa anche verificare.
\item[6.3 V] Si trasforma l'istanza, si risolve con l'algoritmo polinomiale di $P_2$, si ritraduce: tutto polinomiale.
\item[6.4 F] Verso sbagliato. Serve $P\in\NP$ e SAT $\red P$ (un $\NP$-completo noto si riduce a $P$).
\item[6.5 F] $B$ è codificato con $\log B$ bit, quindi $O(NB)$ è esponenziale nella dimensione: pseudo-polinomiale.
\item[6.6 V] Con codifica ragionevole (binaria) dei numeri.
\item[6.7 V] Il sottoinsieme è il certificato; la verifica costa $O(N)$.
\item[6.8 F] Nel caso peggiore.
\item[7.1 F] 0, 1 o infinite: se due punti sono ottimi lo è tutto il segmento.
\item[7.2 F] Regione illimitata non basta: serve anche una direzione di recessione $e$ con $c^Te>0$ (per un max). Es.: $\min x_1$ su $\{x\ge0\}$ ha ottimo.
\item[7.3 V] Per i vincoli $\le$ si aggiunge invece la slack.
\item[7.4 F] Per minimizzare ci si muove nella direzione opposta, $-c$.
\item[8.1 V] Con $\lambda=\frac13$: $\frac13(4,0)+\frac23(1,3)=(2,2)$.
\item[8.2 F] Es.\ l'unione dei due semiassi positivi è un cono non convesso.
\item[8.3 F] È convesso ma non è intersezione di un numero finito di semispazi.
\item[8.4 V] Teorema fondamentale della PL.
\item[8.5 V] Regione convessa e obiettivo lineare.
\item[8.6 F] Non ha vertici: contiene rette (ogni punto sta in mezzo a due punti della striscia).
\end{description}

\chapter{Riepilogo: formulario di modellazione}

\begin{center}
\renewcommand{\arraystretch}{1.3}
\small
\begin{tabular}{@{}p{0.42\textwidth}p{0.52\textwidth}@{}}
\toprule
\textbf{Situazione} & \textbf{Modellazione lineare}\\ \midrule
Frazione minima/massima di una miscela & $\sum_{i\in V}a_ix_i\ge\alpha\sum_{i\in F}a_ix_i$\\
$x>0\Rightarrow y=1$ & $x\le My$\\
Lotto minimo: $x\in\{0\}\cup[q,M]$ & $qy\le x\le My$\\
Costo fisso $f$ se $x>0$ & obiettivo $fy+cx$, vincolo $x\le My$\\
Vincolo attivo solo se $y=1$ & $g(x)\le b+M(1-y)$, $M\ge\max g-b$\\
Almeno uno tra due vincoli & $g_1\le b_1+M_1y,\ g_2\le b_2+M_2(1-y)$\\
$a\Rightarrow b$ & $y_a\le y_b$\\
$a\land b\Rightarrow c$ & $y_a+y_b-1\le y_c$\\
$a\lor b\Rightarrow c$ & $y_a+y_b\le2y_c$\\
almeno $m$ di $n\Rightarrow w$ & $\sum y_i-m+1\le(n-m+1)w$\\
$z=y_1\land y_2$ (prodotto di binarie) & $z\le y_1,\ z\le y_2,\ z\ge y_1+y_2-1$\\
$\min\max_h f_h(x)$ & $\min z,\ z\ge f_h(x)\ \forall h$\\
$\max\min_h f_h(x)$ & $\max z,\ z\le f_h(x)\ \forall h$\\
$\min|f(x)|$ & $\min z,\ z\ge f(x),\ z\ge-f(x)$\\
$\min\sum|\delta_i|$, $\delta_i$ libera & $\delta_i=\delta_i^+-\delta_i^-$, $\min\sum(\delta_i^++\delta_i^-)$\\
Variabile libera & $x=x^+-x^-$, $x^\pm\ge0$\\
Valore in $\{d_1,\dots,d_r\}$ & $\sum_ky_k=1$, $c=\sum_kd_ky_k$\\
Magazzino multi-periodo & $s_{t-1}+x_t=d_t+s_t$\\
Assegnamento con apertura & $\sum_jx_{ij}=1$, $x_{ij}\le y_j$\\
Copertura & $\sum_ja_{ij}x_j\ge1\ \forall i$\\
Turni ciclici & variabili = \emph{inizi} del turno\\
Rapporto $\ge\alpha$ tra espressioni & moltiplica per il denominatore (se $>0$)\\
Disuguaglianza stretta $x>0$ & mai: usa $x\ge\varepsilon$ o variabili binarie\\
\bottomrule
\end{tabular}
\end{center}

\section*{Complessità in una pagina}
\begin{itemize}[nosep]
\item \textbf{Istanza}: problema con dati numerici fissati. \textbf{Dimensione}: bit di codifica (i numeri pesano $\log$ del loro valore).
\item \textbf{Complessità di un algoritmo}: n.\ di operazioni elementari nel caso peggiore, in funzione della dimensione; notazione $O$. Polinomiale $O(n^d)$, esponenziale $O(2^n)$.
\item \textbf{Riconoscimento}: ``esiste $x$ ammissibile con $cx\ge k$?''.
\item $\Pc$: decidibile in tempo polinomiale. $\NP$: risposta SÌ verificabile in tempo polinomiale dato un certificato. $\Pc\subseteq\NP$; $\Pc\overset{?}{=}\NP$ aperto.
\item $P_1\red P_2$: istanze di $P_1$ trasformate in tempo polinomiale in istanze di $P_2$; $P_2$ è almeno difficile quanto $P_1$.
\item $\NP$-completo: in $\NP$ e ogni problema di $\NP$ vi si riduce. Per dimostrarlo: $P\in\NP$ + riduzione \emph{da} un $\NP$-completo noto.
\item SAT è $\NP$-completo (Cook 1971). PLI è $\NP$-completa: clausola $\to\sum_{\text{pos}}x_i+\sum_{\text{neg}}(1-x_i)\ge1$. Zaino $\NP$-completo. PL continua $\in\Pc$.
\end{itemize}

\section*{PL in una pagina}
\begin{itemize}[nosep]
\item Forma standard: $\min c^Tx$, $Ax=b$, $x\ge0$ ($b\ge0$). Forma canonica: $\max c^Tx$, $Ax\le b$, $x\ge0$. Si passa dall'una all'altra con: cambio di segno dell'obiettivo, slack/surplus, $x=x^+-x^-$, $x=-x'$, uguaglianza = due disuguaglianze.
\item Regione ammissibile = poliedro (intersezione finita di semispazi chiusi), quindi convessa. Politopo = poliedro limitato.
\item Esiti: inammissibile, illimitato, ottimo unico, infiniti ottimi. Mai esattamente 2.
\item Vertice: non è combinazione convessa stretta di due punti distinti di $P$ ($\Leftrightarrow$ $n$ vincoli attivi linearmente indipendenti). Direzione di recessione: $Ad\le0$, $d\ne0$.
\item $P=\mathrm{conv}\{v^i\}+\mathrm{cono}\{e^j\}$. Se l'ottimo è finito: $c^Te^j\le0$ per ogni $j$ ed esiste un vertice ottimo (teorema fondamentale).
\item Con $c\ne0$ l'ottimo sta sulla frontiera; due ottimi $\Rightarrow$ infiniti; ottimo locale = globale.
\end{itemize}


\end{document}
